\section{The two-body kernels of $Q_1$, $S_1$ and $K_1$}

For $x$ in a homogeneous basis of $\mathfrak{psl}(4|4)$, $\lambda_x$ is its $\mathrm{ad}_\Delta$ eigenvalue, $-1$ on $\Sigma$, $+1$ on $\Sigma^*$ and $+2$ for $K_b$, and the degree of $x$ in the position is $0$, $1$ and $2$ (Section 10.9). The two-body kernel of a pair of letters along $x$ is the kernel of Proposition \ref{prop:c10-arity1}; this section computes it along the supercharges, the special supercharges and the special conformal fields. Along $\Sigma$ the vertex $S^{(1)}(\xi)$ meets the gauge-fixing operator in a new way: the commutator $[\xi.,Q^{GF}]$ joins the gaugino to the ghost, and $S^{(1)}(\xi)$ sits on a cycle as a line from $\lambda$ to $c$ whose image under $Q$ is nonzero, in contrast to the composite lines along $\Delta$ and $P_b$. The section shows that the component along $x$ of every local functional is nevertheless a sum over its pairs of letters of local two-body kernels, the finite parts of the cycles with one vertex $S^{(1)}(x)$ and one special edge, plus a one-body term; writes these kernels as simplex integrals, independent of the endpoint along $K_b$; computes them on the scalar bilinears, where the two-body part of $Q_1$ on the traceless part is $-2\log\tfrac32$ units times the classical transformation, and derives $S_1$ and $K_1$ there; and evaluates the pair entries of the twisted differential along a rank-one supercharge. Edges carry $P(\varepsilon,\infty)$ along $\Sigma$ and $\Sigma^*$ (Proposition \ref{prop:c8-Tsafe}) and $P(\varepsilon,T)$ at a finite $T$ along $K_b$ (Definition \ref{def:c8-K}). In this section $\psi$ denotes the finite part of the operators of Definition \ref{def:c8-renorm}, and $L:=\log\tfrac32$.

\begin{definition}[The terms of $S^{(1)}(x)$]
\label{def:c11-terms}
For $\xi\in S_+\otimes W$ the functional $S^{(1)}(\xi)$ of \eqref{eq:c2-S1Q} is the sum of the seven terms of Table \ref{tab:c11-terms}: the quadratic terms $T_1,\dots,T_4$ and the three cubic terms $T_5$, with $F_A=dA+\tfrac12[A\wedge A]$ reduced to $V$ and $dA|_V$, $dA|_{W_{\mathbb R}^\perp}$ the parts of $dA$ in $\wedge^2V$ and in $V\wedge W_{\mathbb R}^\perp$; $T_3$ and $T_4$ carry one derivative on the leg $A$. For $\xi\in S_-\otimes W^*$ the table holds with $\lambda|_{S_+\otimes W}$ and $\lambda|_{S_-\otimes W^*}$ exchanged and $\lambda^*|_{S_+\otimes W^*}$ and $\lambda^*|_{S_-\otimes W}$ exchanged. The quadratic part of $S^{(1)}(\xi)$ is $S^V_2(\xi)+S^\perp_2(\xi)$, with $S^V_2(\xi):=T_1+T_3$, the part with the legs $(\lambda,A^*|_V)$ and $(A|_V,\lambda^*)$, and $S^\perp_2(\xi):=T_2+T_4$, the part with the legs $(\lambda,A^*|_{W_{\mathbb R}^\perp})$ and $(A|_{W_{\mathbb R}^\perp},\lambda^*)$. For $\eta\in\Sigma^*$ and an orthonormal basis $e_\mu$ of $V$, \eqref{eq:c2-S1S} reads
\begin{equation}
S^{(1)}(\eta)=\sum_\mu\int x^\mu s_\mu+\int s_0,\qquad \int s_\mu:=S^{(1)}\big(\rho(e_\mu)\eta\big),\qquad s_0:=-\big(\rho(A|_{W_{\mathbb R}^\perp})\eta,\lambda^*\big);
\label{eq:c11-Seta}
\end{equation}
for $\eta\in S_+\otimes W^*$ the spinors $\rho(e_\mu)\eta$ lie in $S_-\otimes W^*$ and $s_0$ has the legs $A|_{W_{\mathbb R}^\perp}$, $\lambda^*|_{S_+\otimes W^*}$, and for $\eta\in S_-\otimes W$ they lie in $S_+\otimes W$ and $s_0$ has the legs $A|_{W_{\mathbb R}^\perp}$, $\lambda^*|_{S_-\otimes W}$. For an even generator with vector field $v$, $S^{(1)}(x)=\int\big((v.A,A^*)-(v.\lambda,\lambda^*)-(v.c,c^*)\big)$ (Definition \ref{def:c7-vdot}).
\end{definition}

\begin{tabular}{l|l|l}
term & functional, $\int\mathrm{dvol}$ understood & legs\\
\hline
$T_1$ & $-\big(\Gamma(\xi,\lambda)|_V,A^*|_V\big)$ & $\lambda|_{S_-\otimes W^*}$, $A^*|_V$\\
$T_2$ & $-\big(\Gamma(\xi,\lambda)|_{W_{\mathbb R}^\perp},A^*|_{W_{\mathbb R}^\perp}\big)$ & $\lambda|_{S_+\otimes W}$, $A^*|_{W_{\mathbb R}^\perp}$\\
$T_3$ & $\tfrac12\big(\rho(dA|_V)\xi,\lambda^*\big)$ & $A|_V$, $\lambda^*|_{S_+\otimes W^*}$\\
$T_4$ & $\tfrac12\big(\rho(dA|_{W_{\mathbb R}^\perp})\xi,\lambda^*\big)$ & $A|_{W_{\mathbb R}^\perp}$, $\lambda^*|_{S_-\otimes W}$\\
$T_5$ & $\tfrac12\big(\rho(\tfrac12[A\wedge A])\xi,\lambda^*\big)$ & $(A|_V,A|_V,\lambda^*|_{S_+\otimes W^*})$, $(A|_V,A|_{W_{\mathbb R}^\perp},\lambda^*|_{S_-\otimes W})$, $(A|_{W_{\mathbb R}^\perp},A|_{W_{\mathbb R}^\perp},\lambda^*|_{S_+\otimes W^*})$\\
\end{tabular}
\tabcaption{The terms of $S^{(1)}(\xi)$ for $\xi\in S_+\otimes W$ and their legs; their two-component forms are in Remark \ref{rem:c11-twocomp}.\label{tab:c11-terms}}

\begin{lemma}[The generators on the regulator and their lines]
\label{lem:c11-lines}
For an odd generator $x$, $x.$ is the odd operator on $\mathscr E$ that induces the bracket with the quadratic part of $S^{(1)}(x)$ (Proposition \ref{prop:c7-susy}); for an even one it is $v.$ of Definition \ref{def:c7-vdot}. Let $g(H):=\int_\varepsilon^Te^{-tH}dt$, so that $P\star=Q^{GF}g(H)$ and $QP\star+P\star Q=e^{-\varepsilon H}-e^{-TH}=:E(H)$ (Proposition \ref{prop:c9-composite} (iii)).
\begin{itemize}
\item (i) For $\xi\in\Sigma$, $\xi.$ maps the slot of $\lambda$ to the slot of $A$ by $\lambda\mapsto+\Gamma(\xi,\lambda)$, the slot of $A$ to the slot of $\lambda$ by $A\mapsto+\tfrac12\rho(dA)\xi$, the slots of $A^*$ and of $\lambda^*$ to each other by the transposed arrows, and vanishes on the slots of $c$ and $c^*$. In particular
\[
\Big\{S^{(1)}(\xi),\int\omega\,(\iota_{w_1}A,\iota_{w_2}A)\Big\}=\int\omega\,\Big[\big(\iota_{w_1}A,(\rho(w_2)\xi,\lambda)\big)+\big(\iota_{w_2}A,(\rho(w_1)\xi,\lambda)\big)\Big].
\]
\item (ii) $[\xi.,Q]=0$, $[\xi.,H]=0$, and $C_\xi:=[\xi.,Q^{GF}]$ is an even operator of degree $-1$ with constant coefficients and of order one, with four components, each nonzero for $\xi\neq0$: $\lambda\mapsto-d^*\big(\Gamma(\xi,\lambda)|_V\big)$ in the slot of $c$, on a plane wave $e^{i(p,x)}$ the map $\lambda\mapsto+i\big(p,\Gamma(\xi,\lambda)|_V\big)$; a component from $A^*|_V$ to $\lambda$, a nonzero multiple of $A^*|_V\mapsto(p,A^*|_V)\,\xi$; and their transposes from $\lambda^*$ to $A|_V$ and from $c^*$ to $\lambda^*$. For $\xi\in S_+\otimes W$ these are $\lambda|_{S_-\otimes W^*}\to c$, $A^*|_V\to\lambda|_{S_+\otimes W}$, $\lambda^*|_{S_+\otimes W^*}\to A|_V$ and $c^*\to\lambda^*|_{S_-\otimes W}$. The variation kernel $(\xi.\otimes1+1\otimes\xi.)P(\varepsilon,T)=(C_\xi\,g(H)\otimes1)K_0$ has its components in $\lambda\otimes A|_V$, $c\otimes\lambda^*$ and their transposes.
\item (iii) The term $T_1$ with both legs on edges $P$, its leg $\lambda$ on the component $\lambda\otimes\lambda$ and its leg $A^*|_V$ on the component $A^*|_V\otimes c$, is the composite line along $\xi$, with the kernel $2(P\star\,\xi.\,P\star\otimes1)K_0$, whose components lie in $c\otimes\lambda$ and $\lambda\otimes c$. For $\xi\in S_+\otimes W$ and $T=\infty$ it is $-2\,\xi\,e^{-2\varepsilon|p|^2}|p|^{-2}(-,-)^{-1}$ on $c\otimes\lambda|_{S_+\otimes W}$ and $+2\,\xi\,e^{-2\varepsilon|p|^2}|p|^{-2}(-,-)^{-1}$ on $\lambda|_{S_+\otimes W}\otimes c$, and
\[
P\star\,\xi.\,P\star=C_\xi\,Q^{GF}g(H)^2-\xi.\,(Q^{GF})^2g(H)^2,\qquad \big[Q,P\star\,\xi.\,P\star\big]=E(H)\,C_\xi\,g(H)\neq0 .
\]
The terms of $S^{(1)}(\xi)$ with a leg $A^*|_{W_{\mathbb R}^\perp}$ or $\lambda^*$ enter a cycle with that leg on the special edge, and a term $T_5$ also with its leg $\lambda^*$ as an output.
\item (iv) For $\eta\in\Sigma^*$, with $\xi_\mu:=\rho(e_\mu)\eta$,
\[
\eta.=\sum_\mu x^\mu\,\xi_\mu.+\eta^0.+Y_\eta,\qquad [\eta.,Q]=0,\qquad [\eta.,H]=2\sum_\mu\xi_\mu.\,\partial_\mu,
\]
where $\eta^0.$ is the constant operator of order zero induced by $s_0$, with the arrows $A|_{W_{\mathbb R}^\perp}\to\lambda$ and $\lambda^*\to A^*|_{W_{\mathbb R}^\perp}$, and $Y_\eta$ is the constant operator of order zero from $\lambda^*$ to $A^*$ produced when the derivative of $T_3$, $T_4$ in $\sum_\mu\int x^\mu\big(S^V_2(\xi_\mu)+S^\perp_2(\xi_\mu)\big)$ is moved off $A$, $\partial_\nu(x^\mu\,\cdot\,)=x^\mu\partial_\nu+\delta^\mu_\nu$. On a cycle $S^{(1)}(\eta)$ forms the line of (iii) at the spinor $\rho(x)\eta$, linear in the position, and $s_0$ enters with its leg $\lambda^*$ on the special edge.
\item (v) For $K_b$, with $v(x)=(x,x)b-2(b,x)x$ and $\operatorname{div}v=-8(b,x)$, on a slot of Weyl weight $w$ with the rotation generators $\Sigma^{\rho\nu}$, and with $\sigma^\rho(D):=[D,x^\rho]$ for a constant-coefficient operator $D$,
\[
[v.,Q^{GF}]=4(b,x)\,Q^{GF}+C_K,\qquad [v.,H]=4(b,x)\,H-4(w-1)(b,\partial)+4b_\nu\Sigma^{\rho\nu}\partial_\rho,\qquad C_K:=2\,\sigma^\rho(Q^{GF})\big(w\,b_\rho-b_\nu\Sigma^{\rho\nu}\big),
\]
$C_K$ of order zero with constant coefficients, zero on $A^*|_{W_{\mathbb R}^\perp}$, with $C_KA|_V=-4(b,A|_V)$ in the slot of $c$. $S^{(1)}(K_b)$ forms on a cycle the composite line from $c$ to $A|_V$ with the kernel $P\star\,v.\,P\star=v.(P\star)^2+[P\star,v.]P\star$, quadratic in the position, at a finite $T$ (Definition \ref{def:c8-K}).
\end{itemize}
In each case every letter, $A^*|_{W_{\mathbb R}^\perp}$, $\lambda^*$ and $c^*$ included, is active (Lemma \ref{lem:c10-active1}).
\end{lemma}

\begin{proof}
(i) The terms $-(\Gamma(\xi,\lambda),A^*)$ and $\tfrac12(\rho(dA)\xi,\lambda^*)$ of \eqref{eq:c2-S1Q} pair the slots $(\lambda,A^*)$ and $(A,\lambda^*)$; the bracket of a quadratic functional with a coordinate functional is the coordinate functional composed with the corresponding arrow (proof of Lemma \ref{lem:c9-SDelta} (i)), and each pairing term gives one arrow between field slots and the transposed arrow between antifield slots (proof of Proposition \ref{prop:c7-susy}). In the reading of Proposition \ref{prop:c1-Q} the term $(dc,A^*)$ of \eqref{eq:c1-free} gives $c\mapsto-dc$; the term $-(\Gamma(\xi,\lambda),A^*)$, whose letter $\lambda$ is odd like $c$, gives $\lambda\mapsto+\Gamma(\xi,\lambda)$, and $[\xi.,Q]=0$ on the slot of $A$, the part of polynomial degree two of $\{S,S^{(1)}(\xi)\}=0$, fixes $A\mapsto+\tfrac12\rho(dA)\xi$. With $\iota_w\Gamma(\xi,\lambda)=(\rho(w)\xi,\lambda)$ [ESW, (28)] the classical component on the bilinear is the first functional of Corollary \ref{cor:c7-Q}.

(ii) $[\xi.,Q]=0$ and $[\xi.,H]=0$ are Proposition \ref{prop:c7-susy}. $Q^{GF}$ maps $A^*\mapsto-2A^*$ into the slot of $A$, $A\mapsto-d^*(A|_V)$ into the slot of $c$, $c^*\mapsto dc^*$ into the slot of $A^*|_V$ and $\lambda^*\mapsto-{\not d}\lambda^*$ into the slot of $\lambda$, and vanishes on $c$ and $\lambda$ (Chapter 3). Hence $\xi.Q^{GF}$ has the arrows $A^*\to A\to\lambda$, $c^*\to A^*|_V\to\lambda^*$, $\lambda^*\to\lambda\to A$, and $Q^{GF}\xi.$ the arrows $\lambda\to A\to c$, $A^*\to\lambda^*\to\lambda$, $\lambda^*\to A^*\to A$; each composite has order one. On $\lambda$, $C_\xi\lambda=Q^{GF}\Gamma(\xi,\lambda)=-d^*\big(\Gamma(\xi,\lambda)|_V\big)$, and $d^*=-\iota_\partial$ is $-i\iota_p$ on plane waves. On $A^*=a\,e^{i(p,x)}$ the two composites to $\lambda$ are opposite multiples of $\rho(p)\rho(a)\xi$ plus, through $\rho(q(p\wedge a))=\rho(p)\rho(a)-(p,a)$ [ESW, (29)] in the first, a multiple of $(p,a)\xi$; for $a\in W_{\mathbb R}^\perp$ the two cancel. The components from $\lambda^*$ and $c^*$ are the transposes, since $\xi.$ preserves the pairing and $Q^{GF}$ has the adjointness of Proposition \ref{prop:c3-adjoint}, so $(C_\xi\otimes1)K_0$ is graded symmetric. $\Gamma(\xi,-)|_V$ is nonzero on one chirality of $\lambda$ for $\xi\neq0$, and the component from $A^*$ to $\lambda$ is nonzero by Proposition \ref{prop:c7-susy}. The variation kernel is the third item of Proposition \ref{prop:c7-susy} with $T$ in place of $\infty$.

(iii) As in Proposition \ref{prop:c9-composite} (i), with the signs of Lemma \ref{lem:c9-ops} for the odd $\xi.$, both orders of the two legs give $(P\star\,\xi.\,P\star\otimes1)K_0$. Since $[\xi.,g(H)]=0$, $[\xi.,P\star]=C_\xi g(H)$, which gives the first identity; with $[\xi.,Q]=0$,
\[
Q\,(P\star\xi.P\star)=E\,\xi.P\star-P\star Q\,\xi.P\star=E\,\xi.P\star+P\star\xi.E-P\star\xi.P\star Q,
\]
so $[Q,P\star\xi.P\star]=E\,(\xi.P\star+P\star\xi.)=E\,C_\xi g$. The components follow the arrows $\lambda^*\to\lambda\to A\to c$ and $c^*\to A^*|_V\to\lambda^*\to\lambda$. On a plane wave the composite on $c^*$ is a multiple of $\rho(p)\rho(p)\xi=|p|^2\xi$ times $g(H)^2=e^{-2\varepsilon|p|^2}|p|^{-4}$ at $T=\infty$, a scalar line of momentum degree $-2$; the multiple, $-2$ on $c\otimes\lambda$, is the value of both evaluations of Remark \ref{rem:c11-oddmethod}. A leg $A^*|_{W_{\mathbb R}^\perp}$ or $\lambda^*$ has no component of $P$ (Proposition \ref{prop:c3-P}).

(iv) By \eqref{eq:c11-Seta} the field arrows of $\int x^\mu S_2(\xi_\mu)$, $S_2$ the quadratic part, are $x^\mu$ times those of $\xi_\mu.$; its antifield arrows, the transposes, are obtained by integrating by parts the derivative of $T_3$, $T_4$ against $x^\mu\lambda^*$, which gives $x^\mu$ times the antifield arrows of $\xi_\mu.$ plus $Y_\eta$. For the odd $Q$, $[x^\mu\xi_\mu.,Q]=x^\mu[\xi_\mu.,Q]+[Q,x^\mu]\xi_\mu.$, and the vanishing of $[\eta.,Q]$ is the part of polynomial degree two of $\{S,S^{(1)}(\eta)\}=0$ (Lemma \ref{lem:c2-components}); it requires $Y_\eta$. $H$ is $-\partial^2$ on every component, $[x^\mu,-\partial^2]=2\partial_\mu$, and $\eta^0.$, $Y_\eta$ and $\xi_\mu.$ have constant coefficients.

(v) For a conformal Killing field, $\partial_\mu v_\nu=\tfrac14(\operatorname{div}v)\delta_{\mu\nu}+\omega_{\mu\nu}$ with $\omega$ antisymmetric and $v.=v^\mu\partial_\mu+\tfrac12\omega_{\mu\nu}\Sigma^{\mu\nu}+\tfrac w4\operatorname{div}v$; the transport, spin and weight terms and the rotation covariance of a constant-coefficient operator $D$ of weight $-2$ give
\[
[v.,D]=-\tfrac12(\operatorname{div}v)D-\tfrac w4\sigma^\rho(D)\,\partial_\rho(\operatorname{div}v)-\tfrac12\sigma^\rho(D)(\partial_\rho\omega_{\mu\nu})\Sigma^{\mu\nu},\qquad [v.,H]=-\tfrac12(\operatorname{div}v)H+\tfrac{w-1}2(\partial_\rho\operatorname{div}v)\partial_\rho+(\partial_\rho\omega_{\mu\nu})\Sigma^{\mu\nu}\partial_\rho,
\]
with $\partial^2v_\mu=-\tfrac12\partial_\mu\operatorname{div}v$. For $K_b$, $\omega_{\mu\nu}=2(x_\mu b_\nu-x_\nu b_\mu)$ and $\partial_\rho\omega_{\mu\nu}=2(\delta_{\rho\mu}b_\nu-\delta_{\rho\nu}b_\mu)$. On $A^*|_{W_{\mathbb R}^\perp}$, $Q^{GF}=-2$ has order zero; on $A|_V$, $\sigma^\rho(Q^{GF})Y=Y_\rho$, $w=1$ and $\sigma^\rho(Q^{GF})\Sigma^{\rho\nu}A=3A_\nu$, so $C_KA|_V=2(b,A|_V)-6(b,A|_V)$. The line is Lemma \ref{lem:c11-kernels} (iii), (iv) with this $v$.
\end{proof}

\begin{definition}[Defect cycles and two-body kernels]
\label{def:c11-defect}
Let $x$ be a generator, $X_a$, $X_b$ two letters of any slots and $0<\varepsilon<T\le\infty$.
\begin{itemize}
\item (i) A defect cycle of $(X_a,X_b)$ along $x$ is a chain $X_a-e_0-V_1-e_1-\cdots-V_n-e_n-X_b$ of Definition \ref{def:c10-chain} whose vertices are vertices of $I$ and exactly one vertex $V_s$ given by a term of $S^{(1)}(x)$, entered and left through two of its legs, the third leg of a cubic term being an output, and whose edges are components of $P(\varepsilon,T)$ and exactly one special edge $e_\sigma=K_\varepsilon-K_0$, with the components of $K_0$. The vertex $V_s$ is placed on a line, with both legs on edges $P$; at the special edge, with one leg on $e_\sigma$; or cubic, a term $T_5$ with two legs on the cycle.
\item (ii) $\mathcal D_x(X_a,X_b;\varepsilon,T)$ is the sum of the defect cycles of the pair, amputated, with the outputs left open as in Definition \ref{def:c10-pair}: the $\delta$-linear graphs of arity one along $x$ of the flow of the defect vertex of Lemma \ref{lem:c8-transfer}, applied to $e^{(J+\delta M(\varepsilon))/\hbar}$ with $M=\int\omega\,X_aX_b$, whose cycle is a defect cycle of the pair. A chain read from $X_a$ carries, beyond the signs of Lemma \ref{lem:c10-sign}, the factor
\[
-(-1)^{n(<\sigma)+|S^{(1)}(x)|\,(n(<s)+[s<\sigma])},
\]
with $n(<j)$ the number of odd outputs at the vertices before the position $j$, $[s<\sigma]=1$ if $V_s$ precedes $e_\sigma$ and $0$ otherwise, and $|S^{(1)}(x)|$ the total parity of $S^{(1)}(x)$, $0$ for odd $x$ and $1$ for even $x$; the leading sign is that of the special edge (Proposition \ref{prop:c11-master}), and for even $x$ the factor is that of Definition \ref{def:c11-sp} times this sign.
\item (iii) Let $\mathrm{ct}_m(M)$ and $\mathrm{lc}(M)$ be the coefficients of $\varepsilon^{-m}$ and of $\log\varepsilon$ in $\partial_\delta I^{CT}[\varepsilon](M)$, and $\mathrm{ct}^{(1)}_{x,m}(M)$ and $\mathrm{lc}^{(1)}_x(M)$ those of the arity-one counterterm of $M$ along $x$, the $\hbar\,\mathrm{Sing}_\varepsilon$ of the cycles through two letters of $M$ with one vertex $S^{(1)}(x)$ on the cycle and no special edge (Definition \ref{def:c11-graphs} for $x\in\{\Delta,P_b\}$). The finite special vertex of the pair is
\[
\mathcal F_x(X_a,X_b):=-\sum_{m\ge1}\frac{(-1)^m}{m!}\Big(\big\{I,\mathrm{ct}^{(1)}_{x,m}(X_aX_b)\big\}_{H^m}+\big\{S^{(1)}(x),\mathrm{ct}_m(X_aX_b)\big\}_{H^m}\Big),
\]
with $\{-,-\}_{H^m}$ the bracket of \eqref{eq:c10-finitesp}.
\item (iv) The two-body kernel of $(X_a,X_b)$ along $x$ is
\[
\Upsilon_x(X_a,X_b):=[\varepsilon^0]\,\mathcal D_x(X_a,X_b;\varepsilon,T)+\mathcal F_x(X_a,X_b),
\]
with $[\varepsilon^0]$ the coefficient of $\varepsilon^0$ in the expansion of Theorem \ref{thm:c11-oddsimplex}: a local density of the outputs at the point of the pair, with values in the tensor product of the fibres and colour spaces of the two letters.
\end{itemize}
For $x\in\{\Delta,P_b\}$ the defect cycles with $V_s$ on a line are the special-edge composite cycles of Definition \ref{def:c11-sp}.
\end{definition}

\begin{theorem}[Two-body form of the components along every generator]
\label{thm:c11-classes}
Let $x$ be a generator and $O=\int\omega\,X_1\cdots X_n$ a monomial local functional, gauge-variant ones included. Let $J=I_{\mathfrak{psl}(4|4)}-I^{CT}_{\mathfrak{psl}(4|4)}[\varepsilon]-\hbar F_1$ with $\hbar F_1$ a one-loop quantisation through arity one along $x$, that is, \eqref{eq:c8-qme} modulo $\hbar^2$ in arity zero and in arity one along $x$, and let $\psi=0$. The endpoint is $T=\infty$ along $\Sigma$, $\Sigma^*$, $\Delta$, $P_b$ and the rotations, and any $T\in(0,\infty)$ along $K_b$.
\begin{itemize}
\item (i) The $\delta$-linear graphs of arity one along $x$ and of order $\hbar$ of the flow of the defect vertex of Lemma \ref{lem:c8-transfer}, applied to $e^{(J+\delta O(\varepsilon))/\hbar}$, form the classes: (F) the one-loop graphs of the vertex $\delta\{S^{(1)}(x),O\}$ and of the vertex $\delta(Q+\{I,-\})O$ with a separate vertex $S^{(1)}(x)$; (C2) the defect cycles of the pairs of letters of $O$, with trees of $I$ on the outputs; (C0) the cycles through two letters of $O$ with the special edge or a contraction through $K_\varepsilon$, with $S^{(1)}(x)$ on a tree; (Br) the graphs with the special edge on a tree; (Ac) the graphs whose cycle passes through vertices of the action alone; (Ct) the trees with one counterterm vertex of order $\hbar$.
\item (ii) (F) is the one-loop part of the arity-one flow of $sO$, and (C0) with the arity-zero counterterm vertices of (Ct) tends to the trees of $\hbar L_1(O)^{(0)}$ with one vertex $S^{(1)}(x)$, both on the left side of \eqref{eq:c8-rep}; (Br) tends to $0$; (Ac) with the defect vertices of the action tends to $0$ by the master equation in arity zero and in arity one along $x$.
\item (iii) With the Koszul signs of \eqref{eq:c10-decomp},
\begin{equation}
\hbar L_1(O)^{(1)}(x)=\sum_{1\le a<b\le n}(-1)^{|X_b|(|X_{a+1}|+\cdots+|X_{b-1}|)}\int\omega\,X_1\cdots X_{a-1}\,\Upsilon_x(X_a,X_b)\,X_{a+1}\cdots\widehat{X_b}\cdots X_n-\big\{\hbar F_1^{(1)}(x),O\big\},
\label{eq:c11-decompx}
\end{equation}
the two-body kernels of Definition \ref{def:c11-defect} summed over the unordered pairs of letters, and a one-body term, the bracket with the arity-one finite counterterm along $x$.
\item (iv) The coefficients of $\log\varepsilon$ and of $\varepsilon^{-m}$ cancel. The defect cycles carry no $\log\varepsilon$ (Theorem \ref{thm:c11-oddsimplex}), and the first cancellation is the logarithmic covariance
\begin{equation}
\mathrm{lc}\big(\{S^{(1)}(x),O\}\big)-\big\{S^{(1)}(x),\mathrm{lc}(O)\big\}=s\,\mathrm{lc}^{(1)}_x(O)-\mathrm{lc}^{(1)}_x(sO)+\big\{L^{(1)}(x),O\big\},
\label{eq:c11-logcov}
\end{equation}
with $L^{(1)}(x)$ the coefficient of $\log\varepsilon$ in $I^{CT}_{\mathfrak{psl}(4|4)}[\varepsilon]^{(1)}(x)$, which has no term $\varepsilon^{-m}$ along $\Sigma$, $\Sigma^*$, $\Delta$, $P_b$ and $K_b$. The second is the pole identity, for $m\ge1$,
\begin{equation}
\begin{aligned}
&\sum_{a<b}\pm\Big([\varepsilon^{-m}]\,\mathcal D_x(X_a,X_b)-\sum_{j\ge1}\tfrac{(-1)^j}{j!}\Big(\big\{I,\mathrm{ct}^{(1)}_{x,m+j}(X_aX_b)\big\}_{H^j}+\big\{S^{(1)}(x),\mathrm{ct}_{m+j}(X_aX_b)\big\}_{H^j}\Big)\Big)\\
&\qquad=s\,\mathrm{ct}^{(1)}_{x,m}(O)+\big\{S^{(1)}(x),\mathrm{ct}_m(O)\big\}-\mathrm{ct}_m\big(\{S^{(1)}(x),O\}\big)-\mathrm{ct}^{(1)}_{x,m}(sO),
\end{aligned}
\label{eq:c11-pole}
\end{equation}
whose right side, for a single pair $O=\int\omega\,X_aX_b$, contains the vacuum poles of the pairs of the cubic monomials of $\{S^{(1)}(x),O\}$ and of $sO$ times their third letter.
\item (v) \eqref{eq:c11-decompx} holds along $K_b$ at every $T\in(0,\infty)$ with the same right side. Given the master equation through arity one along $x$, the local representative of $L(O)$ exists through arity one along $x$ for every local functional $O$: (iii) of Definition \ref{def:c10-hyp} holds along every generator.
\end{itemize}
\end{theorem}

\begin{proof}
(i) At order $\hbar$ a graph of the flow has one cycle and vertices of order $\hbar^0$, or no cycle and one vertex of order $\hbar$ (Proposition \ref{prop:c3-graphs}); the part of arity one along $x$ and linear in $\delta$ places $x$ on one vertex and $\delta$ on one vertex. Expand the defect vertex in $\delta$, in the arity and in $\hbar$, with $I_{\mathfrak{psl}(4|4)}=I+S^{(1)}+S^{(2)}$ and $O(\varepsilon)=O-\partial_\delta I^{CT}_{\mathfrak{psl}(4|4)}[\varepsilon](O)$. The classical master equation \eqref{eq:c2-mc} leaves at order $\hbar^0$ the vertices $\delta\{S^{(1)}(x),O\}$ and $\delta(Q+\{I,-\})O$ and the special-edge parts of $\{S^{(1)}(x),O\}_\varepsilon$, $\{I,O\}_\varepsilon$, $\{I,S^{(1)}(x)\}_\varepsilon$, $\tfrac12\{I,I\}_\varepsilon$, together with $\delta\hbar\Delta_\varepsilon O$; $\hbar\Delta_\varepsilon S^{(1)}(x)$ vanishes, its quadratic terms pairing a field with the antifield of another field and its cubic terms carrying $\lambda^*$ and no $\lambda$. A cycle through a vertex of the first two kinds is in (F). A cycle with a special edge passes through $O$ along two letters (Theorem \ref{thm:c10-decomp} (i)) or avoids $O$, and the positions of $S^{(1)}(x)$ and of the special edge sort it into (C2), (C0), (Br) or (Ac).

(ii) By Theorem \ref{thm:c8-L} (iv) the left side of \eqref{eq:c8-rep} in arity one along $x$ is $\big(L(O)^{(1)}(x)\big)[T]+\big(L(O)^{(0)}\big)[T]^{(1)}(x)$ with $L(O)=sO+\hbar L_1(O)$; the one-loop part of $(sO)[T]^{(1)}(x)$ consists of the graphs (F) and their counterterm trees. For (C0), the arity-zero analysis of Theorem \ref{thm:c10-decomp}, \eqref{eq:c10-finitesp} and Proposition \ref{prop:c10-arity1} applies with a tree carrying $S^{(1)}(x)$ attached at generic momenta. (Br): a special edge on a tree is $e^{-\varepsilon|p|^2}-1$ at a generic momentum times a renormalised subgraph (Theorem \ref{thm:c8-obstruction} (ii), Proposition \ref{prop:c4-defect1}). (Ac): the cycle and its counterterm vertices form the flow of the defect of the action in arity zero or in arity one along $x$, whose limit is $0$ by Definition \ref{def:c8-ob} and the master equation.

(iii) What remains is (C2) with the arity-one vertices of (Ct), against the tree of $\hbar L_1(O)^{(1)}(x)$ and the counterterm trees of $-\partial_\delta I^{CT}[\varepsilon](\{S^{(1)}(x),O\})$ and $-\mathrm{ct}^{(1)}_x(sO)$ on the left. The vertex $-\{\hbar F_1^{(1)}(x),O\}_\varepsilon$ gives the one-body term, its special-edge part tending to $0$. With $\{f,g\}_\varepsilon-\{f,g\}=\sum_{j\ge1}\frac{(-\varepsilon)^j}{j!}\{f,g\}_{H^j}$, the part of order $\varepsilon^0$ of the special-edge part of $-\varepsilon^{-m}\{f,g_m\}_\varepsilon$ is $-\frac{(-1)^m}{m!}\{f,g_m\}_{H^m}$, which for the vertices $-\{I,\mathrm{ct}^{(1)}_x(O)\}_\varepsilon$ and $-\{S^{(1)}(x),\partial_\delta I^{CT}[\varepsilon](O)\}_\varepsilon$ is the finite special vertex; the terms with $j<m$ are poles, those with $j>m$ tend to $0$, and $\log\varepsilon$ times a special edge tends to $0$. Each defect cycle is a sum of poles, $[\varepsilon^0]\mathcal D_x$ and a vanishing remainder (Theorem \ref{thm:c11-oddsimplex}); the trees on its outputs converge, and Theorem \ref{thm:c8-L} (i) identifies the limit with the tree flow of \eqref{eq:c11-decompx}. The sum over the pairs with the Koszul signs is that of Theorem \ref{thm:c10-decomp} (iii).

(iv) The left side has a limit; collecting the coefficients of $\log\varepsilon$ and of $\varepsilon^{-m}$ of (C2), of the special-edge parts of the (Ct) brackets shifted by $j$, of the brackets without special edge and of the counterterm trees moved from the left side gives \eqref{eq:c11-logcov} and \eqref{eq:c11-pole}. A coefficient of $\varepsilon^{-m}$ in $I^{CT}_{\mathfrak{psl}(4|4)}[\varepsilon]^{(1)}(x)$ is a colour-singlet local functional of degree $-1$ and total Weyl weight $4-\tfrac12\lambda_x-2m$, the position at weight $-1$ and of degree $p$ at most that of $x$ (Proposition \ref{prop:c10-bound1}); a colour singlet of degree $-1$ contains an antifield of weight at least $\tfrac52$ and a second adjoint field of weight at least $1$, or $c^*$ and $c$, so its density has weight at least $\tfrac72$ and the total weight is at least $\tfrac72-p$. For $m\ge1$ the total weight is at most $\tfrac52$, $\tfrac32$, $2$, $3$, $1$ along $\Sigma$, $\Sigma^*$, $\Delta$, $P_b$, $K_b$, below $\tfrac72$, $\tfrac52$, $\tfrac52$, $\tfrac72$, $\tfrac32$. Read for one pair, the cubic terms of $\{S^{(1)}(x),O\}$ and of $sO$ have pairs with the empty leg set of Theorem \ref{thm:c10-decomp} (iv), whose poles stand on the right with the third letter.

(v) The argument uses $sO$ through (F), which cancels between the two sides, and through the counterterm trees. Along $K_b$ at a fixed finite $T$ the edges $P(\varepsilon,T)$ are bounded uniformly in $T$ and decay as $e^{-|y|^2/4T}$ (Definition \ref{def:c8-PT}); the graphs of (Br) and (Ac) carry at fixed $\varepsilon$ the term $\tfrac{3}{8\pi^2}\hbar\log T$ times a tree with the factor $e^{-\varepsilon|p|^2}-1$ (Theorem \ref{thm:c8-K}), which tends to $0$; $\Upsilon_{K_b}$ is independent of $T$ by Theorem \ref{thm:c11-oddsimplex} (iii), and Theorem \ref{thm:c8-K} gives the independence of the local representative. The classes (C2) and (Ct) have local limits, so Theorem \ref{thm:c8-L} (i) gives the local representative.
\end{proof}

\begin{theorem}[Defect cycles as simplex integrals]
\label{thm:c11-oddsimplex}
For a defect cycle with $n$ vertices and the edges $e_0,\dots,e_n$, let $q_j$, the centre, the momentum form and the $m$-th simplex integrand $F_m$ be those of Definition \ref{def:c10-cycle} with $n$ in place of $k$, the numerator containing the function of the term of $S^{(1)}(x)$ at $V_s$ (Table \ref{tab:c11-terms}, Definition \ref{def:c7-vdot}) and the special-edge factor $-|l+q_\sigma|^2$ times the fibre factor of $K_0$; along $\eta$ and $K_b$ a momentum $r$ is injected at $V_s$ and carried by $e_s,\dots,e_n$ to $X_b$ and $\omega$. Let $D$ be the superficial degree of Proposition \ref{prop:c11-oddbound}.
\begin{itemize}
\item (i) At generic output momenta and $T=\infty$, for $x\in\Sigma\cup\Sigma^*\cup\{\Delta,P_b\}\cup\mathfrak{so}(V)$, the value of the cycle is $\sum_{m\ge1}\varepsilon^{-m}a_m+a_0+o(1)$ with
\begin{equation}
16\pi^2a_m=\frac1m\int_{\alpha_\sigma\text{ smallest}}\Big[\big(\min_{f\ne\sigma}\alpha_f\big)^m-\alpha_\sigma^m\Big]F_m\,d\alpha\quad(m\ge1),\qquad 16\pi^2a_0=\int_{\alpha_\sigma\text{ smallest}}F_0\,\log\frac{\min_{f\ne\sigma}\alpha_f}{\alpha_\sigma}\,d\alpha,
\label{eq:c11-oddsimplex}
\end{equation}
and no term in $\log\varepsilon$; $a_m$ is a polynomial of degree $D-2m$ in the output momenta and in $r$, and $a_0=0$ for $D<0$.
\item (ii) Along $\eta\in\Sigma^*$, at the point $x_O$ of the pair, $\Upsilon_\eta=\sum_\mu x_O^\mu\,\Upsilon_{\rho(e_\mu)\eta}+\Upsilon^0_\eta$: the position part along $b$ is the kernel along $\rho(b)\eta\in\Sigma$, and the translation-invariant part $\Upsilon^0_\eta$ is $-i\sum_\mu\partial_{r_\mu}$ at $r=0$ of \eqref{eq:c11-oddsimplex} with $S^{(1)}(\rho(e_\mu)\eta)$ at $V_s$, plus \eqref{eq:c11-oddsimplex} with $s_0$ at $V_s$.
\item (iii) Along $K_b$, with $z=x-x_O$,
\[
v(x)=v(x_O)+\big(-2(b,x_O)z+\omega_{x_O}z\big)+v(z),\qquad \omega_{x_O}z:=2\big((x_O,z)b-(b,z)x_O\big),
\]
and correspondingly
\begin{equation}
\Upsilon_{K_b}=v(x_O)^\mu\,\Upsilon_{P_{e_\mu}}-(b,x_O)\,\Upsilon^0_\Delta+\Upsilon^0_{M_{\omega_{x_O}}}+\Upsilon^0_{K_b},
\label{eq:c11-Kdec}
\end{equation}
with $\Upsilon^0_y$ the translation-invariant part along the generator $y$ centred at $x_O$, the rotation $M_{\omega_{x_O}}$ having the vector field $\omega_{x_O}z$; the monomials $z^\mu$ and $z^\mu z^\nu$ of the Killing data act as $-i\partial_{r_\mu}$ and $-\partial_{r_\mu}\partial_{r_\nu}$ at $r=0$. At every finite $T$ the coefficients of the expansion in $\varepsilon$ are given by \eqref{eq:c11-oddsimplex}, and $T\partial_T$ of every defect cycle is $O(\varepsilon)$: $\Upsilon_{K_b}$ is independent of $T$.
\item (iv) For a cycle with $n_t=n+1$ edges and a constant integrand $F_0$ ($D=0$), $16\pi^2a_0=F_0\log\frac{n_t}{n_t-1}\big/(n_t-1)!$, so that $16\pi^2a_0=\tfrac12F_0\log\tfrac32$ for $n_t=3$ and $\tfrac16F_0\log\tfrac43$ for $n_t=4$; for a constant $F_1$ ($D=2$), $16\pi^2a_1=F_1\big/\big(n_t^2(n_t-1)(n_t-1)!\big)$. For $n=1$, the vertex $V_s$ at the special edge, with $\alpha:=\alpha_\sigma\in[0,\tfrac12]$,
\[
16\pi^2a_0=\int_0^{1/2}F_0(\alpha)\log\frac{1-\alpha}\alpha\,d\alpha,\qquad 16\pi^2a_m=\frac1m\int_0^{1/2}\big((1-\alpha)^m-\alpha^m\big)F_m(\alpha)\,d\alpha,
\]
in $\mathbb Q+\mathbb Q\log2$ for polynomial numerators; for $n=2$ with $V_s$ on a line between a letter or a vertex of $I$ and the special edge, Proposition \ref{prop:c11-k1} applies, with values in $\mathbb Q+\mathbb Q\log2+\mathbb Q\log3$.
\end{itemize}
\end{theorem}

\begin{proof}
(i) The domain $\{t_\sigma\le\varepsilon,\ t_f\ge\varepsilon\ (f\ne\sigma)\}$ is the difference of the domains with the cutoff shapes $1/\min_{f\ne\sigma}\alpha_f$ and $1/\min_f\alpha_f$, which differ on the region where $\alpha_\sigma$ is the smallest coordinate; Proposition \ref{prop:c10-shape} gives for each shape the poles, the same coefficient of $\log\varepsilon$ and the finite part, and the difference is \eqref{eq:c11-oddsimplex}, as in Lemma \ref{lem:c10-special}. The region $\tau\ge1$ converges at $T=\infty$: the edges carry distinct generic momenta (proof of Theorem \ref{thm:c10-decomp} (ii)), the composite line along $\xi$ is of order $|p|^{-2}$ at small momentum (Lemma \ref{lem:c11-lines} (iii)), the composite kernels along $\Delta$, $P_b$ and $\mathfrak{so}(V)$ of order $|p|^{-3}$ (Proposition \ref{prop:c9-composite} (iv)), and along $\eta$ the line carries the spinor $\rho(b)\eta$ in the position part and one derivative in $r$ in the translation-invariant part; the measure of $\{\alpha_\sigma\le\varepsilon/\tau\}$ is of order $\varepsilon$. The coefficient of $\tau^{-1-m}$ after the Gaussian integral is homogeneous of degree $D-2m$. (ii) By \eqref{eq:c11-Seta}, insert $x^\mu=x_O^\mu+(x-x_O)^\mu$ and argue as in Proposition \ref{prop:c11-position} (i): on plane waves the factor $(x-x_O)^\mu$ is $-i\partial/\partial r_\mu$ of the momentum injected at $V_s$ and absorbed by $\omega$, which commutes with the Gaussian integral and the simplex integration; the same holds for the special edges, the counterterm vertices and the finite special vertices. (iii) Expand $(x_O+z,x_O+z)b-2(b,x_O+z)(x_O+z)$; $S^{(1)}$ of Definition \ref{def:c7-vdot} is linear in $v$ with its weight and spin terms, and $\operatorname{div}v(x)=-8(b,x_O)-8(b,z)$, $\operatorname{div}(2z)=8$, $\operatorname{div}(\omega_{x_O}z)=0$. At a finite $T$ the region $\tau\ge\tau_0$ is compact in $\tau$, so the expansion in $\varepsilon$ is that of the region of small $\tau$. $T\partial_T$ replaces one edge $e$ by $TL_T$, of scalar factor $Te^{-T|l+q_e|^2}$, which confines the loop momentum, while the special edge carries $e^{-\varepsilon|l+q_\sigma|^2}-1=O\big(\varepsilon(1+|l|)^2\big)$; the polynomial dependence on the position enters through derivatives in $r$, which preserve these bounds, and the finite special vertices are built from the coefficients $\mathrm{ct}_m$, $\mathrm{ct}^{(1)}_{x,m}$, independent of $T$ (Proposition \ref{prop:c8-Tindep}). (iv) For a constant integrand the integrals of \eqref{eq:c11-oddsimplex} over the region where $\alpha_\sigma$ is the smallest coordinate are elementary; for $n=1$ the simplex is $[0,1]$, the region is $\alpha_\sigma\le\tfrac12$ and $\min_{f\ne\sigma}\alpha_f=1-\alpha_\sigma$, and $\int_0^{1/2}\alpha^a(1-\alpha)^b\log\frac{1-\alpha}\alpha\,d\alpha\in\mathbb Q+\mathbb Q\log2$. Proposition \ref{prop:c11-k1} uses that the two segments adjacent to $V_s$ carry the same momentum at $r=0$ and that the special edge is the third edge.
\end{proof}

\begin{example}[The cubic term at the special edge]
\label{ex:c11-cubic}
The defect cycles with $n=1$ and an output have a cubic term of $S^{(1)}(\xi)$ at the special edge. For $\xi\in S_+\otimes W$ and the pair $(A|_{W_{\mathbb R}^\perp},\lambda|_{S_+\otimes W})$, the cycle
\[
A|_{W_{\mathbb R}^\perp}\xrightarrow{\ P\ }T_5\big(A|_{W_{\mathbb R}^\perp},A|_{W_{\mathbb R}^\perp},\lambda^*|_{S_+\otimes W^*}\big)\xrightarrow{\ K_\varepsilon-K_0\ }\lambda|_{S_+\otimes W},
\]
with the output $A|_{W_{\mathbb R}^\perp}$ of momentum $q$, has by Theorem \ref{thm:c11-oddsimplex} (iv) and Lemma \ref{lem:c10-special} the value
\[
\frac{\hbar}{16\pi^2}\Big(-\frac1{2\varepsilon}-\frac38|q|^2\Big)
\]
times the value of $T_5$ on the fibre elements and colour factors of the letters and of the output, with the fibre factor $2(-,-)^{-1}$ of the edge $A\otimes A$ and the inverse fibre pairing of $(\lambda,\lambda^*)$ on the special edge. Its contribution to $\Upsilon_\xi\big(A|_{W_{\mathbb R}^\perp},\lambda|_{S_+\otimes W}\big)$ on the leg set $A|_{W_{\mathbb R}^\perp}$ is $\tfrac38\frac{\hbar}{16\pi^2}$ times $\partial^2$ of the output times that factor, and its pole enters \eqref{eq:c11-pole} against $\mathrm{ct}_1$ of the pair $(\lambda|_{S_+\otimes W},\lambda|_{S_+\otimes W})$ of the classical transformation $\{S^{(1)}(\xi),A|_{W_{\mathbb R}^\perp}\}=\Gamma(\xi,\lambda)|_{W_{\mathbb R}^\perp}$, whose coefficient of $\varepsilon^{-1}$ has the leg set $A|_{W_{\mathbb R}^\perp}$ (Table \ref{tab:c10-legsets}).
\end{example}

\begin{proposition}[Output bounds along $\Sigma$, $\Sigma^*$ and $K_b$]
\label{prop:c11-oddbound}
A defect cycle of $(X_a,X_b)$ along $x$ whose outputs have total Weyl weight $w$, fields and derivatives, has the superficial degree
\[
D=w_{X_a}+w_{X_b}+\delta-w,
\]
with $\delta=\tfrac12$ along $\xi\in\Sigma$ and for the position part along $\eta$, $\delta=-\tfrac12$ for $\Upsilon^0_\eta$, and $\delta=1,0,-1$ for the parts of $\Upsilon_{K_b}$ multiplying $v(x_O)$, the first derivatives of $v$ at $x_O$ and its second derivatives. Hence:
\begin{itemize}
\item (i) along $\xi\in\Sigma$: $\Upsilon_\xi(X_a,X_b)$ is translation invariant, of Weyl weight $w_{X_a}+w_{X_b}+\tfrac12$, of the pair degree and of parity $|X_a|+|X_b|+1$, with at most $w_{X_a}+w_{X_b}+\tfrac12$ plus the pair degree output fields, as for the second-order operator on pairs of letters of [BK, (3.7)];
\item (ii) along $\eta\in\Sigma^*$: $\Upsilon_\eta=\sum_\mu x^\mu\,\Upsilon_{\rho(e_\mu)\eta}+\Upsilon^0_\eta$, of parity $|X_a|+|X_b|+1$, with $\Upsilon^0_\eta$ translation invariant of Weyl weight $w_{X_a}+w_{X_b}-\tfrac12$ and at most $w_{X_a}+w_{X_b}-\tfrac12$ plus the pair degree output fields;
\item (iii) along $K_b$: $\Upsilon_{K_b}$ is a polynomial of degree at most two in the position, of Weyl weight $w_{X_a}+w_{X_b}-1$ with the position at weight $-1$, of the pair degree and of parity $|X_a|+|X_b|$, with at most $w_{X_a}+w_{X_b}+1$ plus the pair degree output fields; $\Upsilon^0_{K_b}$ has Weyl weight $w_{X_a}+w_{X_b}-1$.
\end{itemize}
The same bounds hold for the counterterm cycles, whose outputs carry the pair degree minus one, and for the finite special vertices, which are supported on the leg sets with $D\ge2$. These are the bounds of Proposition \ref{prop:c10-bound1} with $\lambda_x=-1,+1,+2$, valid without (iii) of Definition \ref{def:c10-hyp} by Theorem \ref{thm:c11-classes} (v).
\end{proposition}

\begin{proof}
Lemma \ref{lem:c10-degree} with one vertex of weight $4+\delta$: every vertex of $I$ has weight $4$, the term of $S^{(1)}(x)$ has weight $4+\delta$ on its fields and derivatives once the Killing data carry the weights of Proposition \ref{prop:c7-select} ($-\tfrac12$ for a constant spinor, $-1$, $0$, $+1$ for the value and the first and second derivatives of a vector field, the position at weight $-1$), every edge $P$ has numerator degree $w_1+w_2-2$, and the special edge, $-|l+q|^2$ over a Schwinger interval $[0,\varepsilon]$, has degree $0=w+(4-w)-4$. The degree of the coefficient of $\varepsilon^0$ is Theorem \ref{thm:c11-oddsimplex} (i). $S^{(1)}(x)$ has degree $-1$ and total parity $0$ for odd $x$ and $1$ for even $x$, the special edge degree $+1$ and parity $1$. Weyl weight plus degree is at least $1$ for every field and antifield, with equality for $c$ and $A$ (Proposition \ref{prop:c10-bound}), which bounds the number of outputs. The position parts are Theorem \ref{thm:c11-oddsimplex} (ii), (iii). On a leg set of degree $D$ the term of $\mathcal F_x$ with $\mathrm{ct}_m$ or $\mathrm{ct}^{(1)}_{x,m}$ contains a coefficient of $\varepsilon^{-m}$, of degree $D-2m\ge0$ in the momenta, and $m\ge1$.
\end{proof}

\begin{proposition}[Census of the defect cycles]
\label{prop:c11-oddcensus}
\begin{itemize}
\item (i) Every defect cycle conserves the degree and the centre $\mathbb Z/4$ of $SL(4,\mathbb C)$, $W\mapsto iW$, once $x$ carries its charge, $+1$ on $S_+\otimes W$ and on $S_-\otimes W$, $-1$ on $S_-\otimes W^*$ and on $S_+\otimes W^*$, $0$ on the even part; it shifts the total parity by that of $S^{(1)}(x)$ plus one.
\item (ii) Along every generator, defect cycles exist for $45$ of the $55$ unordered pairs of the ten slots; the other ten are the pairs of two letters among $A^*|_{W_{\mathbb R}^\perp}$, $\lambda^*|_{S_+\otimes W^*}$, $\lambda^*|_{S_-\otimes W}$, $c^*$, each of which needs the special edge.
\item (iii) A defect cycle without vertices of $I$ has $V_s$ at the special edge, letter--$P$--$V_s$--special edge--letter; a cycle without outputs is homogeneous in $\varepsilon$ of degree $-\tfrac D2$ (Lemma \ref{lem:c10-self}), so its coefficient of $\varepsilon^0$ is nonzero at $D=0$ alone.
\item (iv) Along $\xi\in S_+\otimes W$ the leg sets through the pairs of field and ghost letters without derivatives are those of Table \ref{tab:c11-defectpairs}; the numbers of leg sets with $n\le3$ derivatives on the two letters, along $\xi\in S_+\otimes W$ and for the translation-invariant parts along $\eta\in S_+\otimes W^*$ and along $K_b$, are those of Table \ref{tab:c11-defectcounts}. Along $S_-\otimes W^*$ and $S_-\otimes W$ the tables hold with the chiralities exchanged, and for the position parts with $\delta$ of Proposition \ref{prop:c11-oddbound}.
\item (v) Along $\xi\in S_+\otimes W$ the pair $(c,c)$ has the empty set of leg sets at the pair weight without derivatives, and the pairs of $c$ with $A|_V$, $A|_{W_{\mathbb R}^\perp}$, $\lambda|_{S_+\otimes W}$, $\lambda|_{S_-\otimes W^*}$ have leg sets with one ghost output.
\end{itemize}
\end{proposition}

\begin{proof}
(i) Each vertex of $I$, each term of $S^{(1)}(x)$ and each component of $P$ and of $K_0$ carries these charges (Table \ref{tab:c11-terms}, Definition \ref{def:c1-vertex}, Proposition \ref{prop:c3-P}). (ii)--(v) Exhaustive enumeration of the chains with up to six vertices on the cycle, exactly one term of $S^{(1)}(x)$ in each of its placements and exactly one special edge, with output weights up to $w_{X_a}+w_{X_b}+3+\delta$; every chain is checked for degree, parity and the centre, and its superficial degree against the count of Lemma \ref{lem:c10-degree}. A letter $A^*|_{W_{\mathbb R}^\perp}$, $\lambda^*$ or $c^*$ has no component of $P$ and lies on the special edge, of which a defect cycle has one.
\end{proof}

\begin{tabular}{l|c|l}
pair & output weight at most & leg sets along $\xi\in S_+\otimes W$: fewest vertices of $I$; number of chains\\
\hline
$(A|_V,A|_V)$ & $5/2$ & $\lambda|_{S_-\otimes W^*}$ $(1;2)$, $A|_V\,\lambda|_{S_-\otimes W^*}$ $(1;8)$, $A|_{W_{\mathbb R}^\perp}\,\lambda|_{S_+\otimes W}$ $(1;7)$, $c\,\lambda^*|_{S_+\otimes W^*}$ $(1;4)$\\
$(A|_V,A|_{W_{\mathbb R}^\perp})$ & $5/2$ & $\lambda|_{S_+\otimes W}$ $(1;4)$, $A|_V\,\lambda|_{S_+\otimes W}$ $(1;15)$, $A|_{W_{\mathbb R}^\perp}\,\lambda|_{S_-\otimes W^*}$ $(1;15)$, $c\,\lambda^*|_{S_-\otimes W}$ $(1;6)$\\
$(A|_V,\lambda|_{S_+\otimes W})$ & $3$ & $\emptyset$ $(0;2)$, $A|_V$ $(0;6)$, $A|_V^2$ $(1;13)$, $(A|_{W_{\mathbb R}^\perp})^2$ $(1;11)$, $A|_V^3$ $(1;25)$, $A|_V\,(A|_{W_{\mathbb R}^\perp})^2$ $(1;55)$, $c\,A^*|_V$ $(2;9)$, $\lambda|_{S_+\otimes W}\,\lambda|_{S_-\otimes W^*}$ $(2;15)$\\
$(A|_V,\lambda|_{S_-\otimes W^*})$ & $3$ & $A|_{W_{\mathbb R}^\perp}$ $(0;5)$, $A|_V\,A|_{W_{\mathbb R}^\perp}$ $(1;19)$, $A|_V^2\,A|_{W_{\mathbb R}^\perp}$ $(1;51)$, $(A|_{W_{\mathbb R}^\perp})^3$ $(1;19)$, $c\,A^*|_{W_{\mathbb R}^\perp}$ $(2;2)$, $(\lambda|_{S_-\otimes W^*})^2$ $(2;7)$, $(\lambda|_{S_+\otimes W})^2$ $(2;6)$\\
$(A|_V,c)$ & $3/2$ & $\lambda|_{S_-\otimes W^*}\,c$ $(2;10)$\\
$(A|_{W_{\mathbb R}^\perp},A|_{W_{\mathbb R}^\perp})$ & $5/2$ & $\lambda|_{S_-\otimes W^*}$ $(1;2)$, $A|_V\,\lambda|_{S_-\otimes W^*}$ $(1;7)$, $A|_{W_{\mathbb R}^\perp}\,\lambda|_{S_+\otimes W}$ $(1;8)$, $c\,\lambda^*|_{S_+\otimes W^*}$ $(1;2)$\\
$(A|_{W_{\mathbb R}^\perp},\lambda|_{S_+\otimes W})$ & $3$ & $A|_{W_{\mathbb R}^\perp}$ $(0;6)$, $A|_V\,A|_{W_{\mathbb R}^\perp}$ $(1;22)$, $A|_V^2\,A|_{W_{\mathbb R}^\perp}$ $(1;59)$, $(A|_{W_{\mathbb R}^\perp})^3$ $(1;21)$, $c\,A^*|_{W_{\mathbb R}^\perp}$ $(2;7)$, $(\lambda|_{S_-\otimes W^*})^2$ $(2;7)$, $(\lambda|_{S_+\otimes W})^2$ $(2;8)$\\
$(A|_{W_{\mathbb R}^\perp},\lambda|_{S_-\otimes W^*})$ & $3$ & $\emptyset$ $(0;2)$, $A|_V$ $(0;5)$, $A|_V^2$ $(1;10)$, $(A|_{W_{\mathbb R}^\perp})^2$ $(1;11)$, $A|_V^3$ $(1;18)$, $A|_V\,(A|_{W_{\mathbb R}^\perp})^2$ $(1;52)$, $\lambda|_{S_+\otimes W}\,\lambda|_{S_-\otimes W^*}$ $(2;13)$\\
$(A|_{W_{\mathbb R}^\perp},c)$ & $3/2$ & $\lambda|_{S_+\otimes W}\,c$ $(2;9)$\\
$(\lambda|_{S_+\otimes W},\lambda|_{S_+\otimes W})$ & $7/2$ & $\lambda|_{S_+\otimes W}$ $(1;3)$, $A|_V\,\lambda|_{S_+\otimes W}$ $(1;11)$, $A|_{W_{\mathbb R}^\perp}\,\lambda|_{S_-\otimes W^*}$ $(1;9)$, $c\,\lambda^*|_{S_-\otimes W}$ $(2;7)$, $A|_V^2\,\lambda|_{S_+\otimes W}$ $(2;27)$, $A|_V\,A|_{W_{\mathbb R}^\perp}\,\lambda|_{S_-\otimes W^*}$ $(2;42)$, $A|_V\,c\,\lambda^*|_{S_-\otimes W}$ $(2;30)$, $(A|_{W_{\mathbb R}^\perp})^2\,\lambda|_{S_+\otimes W}$ $(2;21)$, $A|_{W_{\mathbb R}^\perp}\,c\,\lambda^*|_{S_+\otimes W^*}$ $(3;9)$\\
$(\lambda|_{S_+\otimes W},\lambda|_{S_-\otimes W^*})$ & $7/2$ & $\lambda|_{S_-\otimes W^*}$ $(1;5)$, $A|_V\,\lambda|_{S_-\otimes W^*}$ $(1;18)$, $A|_{W_{\mathbb R}^\perp}\,\lambda|_{S_+\otimes W}$ $(1;17)$, $c\,\lambda^*|_{S_+\otimes W^*}$ $(2;7)$, $A|_V^2\,\lambda|_{S_-\otimes W^*}$ $(2;44)$, $A|_V\,A|_{W_{\mathbb R}^\perp}\,\lambda|_{S_+\otimes W}$ $(2;79)$, $A|_V\,c\,\lambda^*|_{S_+\otimes W^*}$ $(2;30)$, $(A|_{W_{\mathbb R}^\perp})^2\,\lambda|_{S_-\otimes W^*}$ $(2;40)$, $A|_{W_{\mathbb R}^\perp}\,c\,\lambda^*|_{S_-\otimes W}$ $(2;21)$\\
$(\lambda|_{S_+\otimes W},c)$ & $2$ & $c$ $(1;5)$, $A|_V\,c$ $(1;17)$, $A|_V^2\,c$ $(2;39)$, $(A|_{W_{\mathbb R}^\perp})^2\,c$ $(2;18)$\\
$(\lambda|_{S_-\otimes W^*},\lambda|_{S_-\otimes W^*})$ & $7/2$ & $\lambda|_{S_+\otimes W}$ $(1;2)$, $A|_V\,\lambda|_{S_+\otimes W}$ $(1;7)$, $A|_{W_{\mathbb R}^\perp}\,\lambda|_{S_-\otimes W^*}$ $(1;8)$, $A|_V^2\,\lambda|_{S_+\otimes W}$ $(2;17)$, $A|_V\,A|_{W_{\mathbb R}^\perp}\,\lambda|_{S_-\otimes W^*}$ $(2;37)$, $(A|_{W_{\mathbb R}^\perp})^2\,\lambda|_{S_+\otimes W}$ $(2;19)$, $A|_{W_{\mathbb R}^\perp}\,c\,\lambda^*|_{S_+\otimes W^*}$ $(2;12)$\\
$(\lambda|_{S_-\otimes W^*},c)$ & $2$ & $A|_{W_{\mathbb R}^\perp}\,c$ $(1;10)$, $A|_V\,A|_{W_{\mathbb R}^\perp}\,c$ $(2;44)$\\
$(c,c)$ & $1/2$ & $-$\\
\end{tabular}
\tabcaption{The leg sets of the defect cycles along $\xi\in S_+\otimes W$ through the pairs of field and ghost letters without derivatives, with the fewest vertices of $I$ and the number of chains; along $\xi\in S_-\otimes W^*$ the chiralities of $\lambda$ and $\lambda^*$ are exchanged. A leg set enters $\Upsilon_\xi$ when its output weight is at most the bound of the second column, and the empty leg set $\emptyset$ enters the poles.\label{tab:c11-defectpairs}}

\begin{tabular}{l|cccc|cccc|cccc}
pair & $\xi$, $n=0$ & $1$ & $2$ & $3$ & $\Upsilon^0_\eta$, $n=0$ & $1$ & $2$ & $3$ & $\Upsilon^0_{K_b}$, $n=0$ & $1$ & $2$ & $3$\\
\hline
$(A|_V,A|_V)$ & 4 & 9 & 20 & 43 & 1 & 4 & 9 & 20 & 2 & 4 & 8 & 19\\
$(A|_V,A|_{W_{\mathbb R}^\perp})$ & 4 & 9 & 20 & 43 & 1 & 4 & 9 & 20 & 1 & 2 & 7 & 16\\
$(A|_V,\lambda|_{S_+\otimes W})$ & 8 & 19 & 35 & 68 & 2 & 7 & 16 & 33 & 1 & 4 & 9 & 20\\
$(A|_V,\lambda|_{S_-\otimes W^*})$ & 7 & 16 & 33 & 61 & 4 & 8 & 19 & 35 & 1 & 4 & 9 & 20\\
$(A|_V,c)$ & 1 & 4 & 9 & 20 & 0 & 1 & 4 & 9 & 1 & 2 & 4 & 8\\
$(A|_V,A^*|_V)$ & 8 & 22 & 49 & 97 & 3 & 8 & 22 & 49 & 1 & 6 & 15 & 35\\
$(A|_V,A^*|_{W_{\mathbb R}^\perp})$ & 8 & 22 & 49 & 97 & 3 & 8 & 22 & 49 & 1 & 5 & 13 & 29\\
$(A|_V,\lambda^*|_{S_+\otimes W^*})$ & 3 & 11 & 26 & 62 & 1 & 6 & 15 & 35 & 1 & 3 & 8 & 21\\
$(A|_V,\lambda^*|_{S_-\otimes W})$ & 6 & 15 & 35 & 73 & 0 & 3 & 11 & 26 & 1 & 3 & 8 & 21\\
$(A|_V,c^*)$ & 3 & 13 & 34 & 80 & 0 & 3 & 13 & 34 & 1 & 3 & 10 & 25\\
$(A|_{W_{\mathbb R}^\perp},A|_{W_{\mathbb R}^\perp})$ & 4 & 9 & 19 & 39 & 1 & 4 & 9 & 19 & 2 & 4 & 7 & 14\\
$(A|_{W_{\mathbb R}^\perp},\lambda|_{S_+\otimes W})$ & 7 & 16 & 33 & 61 & 4 & 7 & 16 & 30 & 1 & 3 & 7 & 16\\
$(A|_{W_{\mathbb R}^\perp},\lambda|_{S_-\otimes W^*})$ & 7 & 16 & 30 & 60 & 2 & 7 & 16 & 33 & 1 & 3 & 7 & 16\\
$(A|_{W_{\mathbb R}^\perp},c)$ & 1 & 4 & 9 & 20 & 0 & 1 & 4 & 9 & 0 & 1 & 2 & 7\\
$(A|_{W_{\mathbb R}^\perp},A^*|_V)$ & 8 & 22 & 49 & 97 & 3 & 8 & 22 & 49 & 1 & 5 & 13 & 29\\
$(A|_{W_{\mathbb R}^\perp},A^*|_{W_{\mathbb R}^\perp})$ & 7 & 18 & 40 & 81 & 3 & 7 & 18 & 40 & 0 & 1 & 5 & 16\\
$(A|_{W_{\mathbb R}^\perp},\lambda^*|_{S_+\otimes W^*})$ & 3 & 10 & 26 & 59 & 1 & 5 & 14 & 31 & 0 & 1 & 4 & 14\\
$(A|_{W_{\mathbb R}^\perp},\lambda^*|_{S_-\otimes W})$ & 5 & 14 & 31 & 70 & 0 & 3 & 10 & 26 & 0 & 1 & 4 & 14\\
$(A|_{W_{\mathbb R}^\perp},c^*)$ & 3 & 13 & 34 & 80 & 0 & 3 & 13 & 34 & 0 & 1 & 4 & 19\\
$(\lambda|_{S_+\otimes W},\lambda|_{S_+\otimes W})$ & 9 & 20 & 42 & 79 & 3 & 7 & 16 & 35 & 2 & 6 & 12 & 24\\
$(\lambda|_{S_+\otimes W},\lambda|_{S_-\otimes W^*})$ & 9 & 20 & 43 & 82 & 4 & 9 & 20 & 43 & 4 & 7 & 15 & 27\\
$(\lambda|_{S_+\otimes W},c)$ & 4 & 8 & 19 & 35 & 1 & 2 & 7 & 16 & 0 & 1 & 4 & 9\\
$(\lambda|_{S_+\otimes W},A^*|_V)$ & 15 & 35 & 73 & 143 & 5 & 14 & 31 & 70 & 3 & 8 & 21 & 46\\
$(\lambda|_{S_+\otimes W},A^*|_{W_{\mathbb R}^\perp})$ & 14 & 31 & 70 & 135 & 3 & 10 & 26 & 59 & 1 & 4 & 14 & 35\\
$(\lambda|_{S_+\otimes W},\lambda^*|_{S_+\otimes W^*})$ & 8 & 21 & 47 & 88 & 3 & 8 & 21 & 47 & 0 & 2 & 7 & 19\\
$(\lambda|_{S_+\otimes W},\lambda^*|_{S_-\otimes W})$ & 8 & 21 & 46 & 86 & 1 & 3 & 11 & 29 & 0 & 1 & 5 & 14\\
$(\lambda|_{S_+\otimes W},c^*)$ & 10 & 25 & 62 & 116 & 2 & 5 & 19 & 49 & 0 & 2 & 10 & 29\\
$(\lambda|_{S_-\otimes W^*},\lambda|_{S_-\otimes W^*})$ & 7 & 16 & 35 & 68 & 4 & 9 & 20 & 42 & 2 & 6 & 12 & 24\\
$(\lambda|_{S_-\otimes W^*},c)$ & 2 & 7 & 16 & 33 & 2 & 4 & 8 & 19 & 0 & 1 & 4 & 9\\
$(\lambda|_{S_-\otimes W^*},A^*|_V)$ & 14 & 31 & 70 & 135 & 6 & 15 & 35 & 73 & 3 & 8 & 21 & 46\\
$(\lambda|_{S_-\otimes W^*},A^*|_{W_{\mathbb R}^\perp})$ & 10 & 26 & 59 & 123 & 5 & 14 & 31 & 70 & 1 & 4 & 14 & 35\\
$(\lambda|_{S_-\otimes W^*},\lambda^*|_{S_+\otimes W^*})$ & 3 & 11 & 29 & 60 & 3 & 8 & 21 & 46 & 0 & 1 & 5 & 14\\
$(\lambda|_{S_-\otimes W^*},\lambda^*|_{S_-\otimes W})$ & 8 & 21 & 47 & 88 & 3 & 8 & 21 & 47 & 0 & 2 & 7 & 19\\
$(\lambda|_{S_-\otimes W^*},c^*)$ & 5 & 19 & 49 & 101 & 3 & 10 & 25 & 62 & 0 & 2 & 10 & 29\\
$(c,c)$ & 0 & 1 & 4 & 9 & 0 & 0 & 1 & 4 & 0 & 1 & 2 & 4\\
$(c,A^*|_V)$ & 9 & 20 & 43 & 82 & 4 & 9 & 20 & 43 & 4 & 8 & 19 & 35\\
$(c,A^*|_{W_{\mathbb R}^\perp})$ & 9 & 20 & 43 & 82 & 4 & 9 & 20 & 43 & 2 & 7 & 16 & 32\\
$(c,\lambda^*|_{S_+\otimes W^*})$ & 6 & 14 & 30 & 53 & 4 & 8 & 19 & 35 & 1 & 4 & 9 & 20\\
$(c,\lambda^*|_{S_-\otimes W})$ & 8 & 19 & 35 & 64 & 2 & 6 & 14 & 30 & 1 & 4 & 9 & 20\\
$(c,c^*)$ & 8 & 22 & 49 & 91 & 3 & 8 & 22 & 49 & 1 & 6 & 15 & 35\\
$(A^*|_V,A^*|_V)$ & 13 & 34 & 80 & 169 & 3 & 13 & 34 & 80 & 3 & 10 & 25 & 61\\
$(A^*|_V,A^*|_{W_{\mathbb R}^\perp})$ & 13 & 34 & 80 & 169 & 3 & 13 & 34 & 80 & 1 & 4 & 19 & 49\\
$(A^*|_V,\lambda^*|_{S_+\otimes W^*})$ & 5 & 19 & 49 & 101 & 3 & 10 & 25 & 62 & 0 & 2 & 10 & 29\\
$(A^*|_V,\lambda^*|_{S_-\otimes W})$ & 10 & 25 & 62 & 116 & 2 & 5 & 19 & 49 & 0 & 2 & 10 & 29\\
$(A^*|_V,c^*)$ & 4 & 15 & 46 & 101 & 1 & 4 & 15 & 46 & 0 & 1 & 7 & 24\\
\end{tabular}
\tabcaption{The number of leg sets of the defect cycles entering $\Upsilon_\xi$, $\xi\in S_+\otimes W$, the translation-invariant part $\Upsilon^0_\eta$, $\eta\in S_+\otimes W^*$, and the translation-invariant part $\Upsilon^0_{K_b}$, for $n$ derivatives on the two letters and up to six vertices on the cycle, on the $45$ pairs of slots with defect cycles; the other ten pairs, of two letters among $A^*|_{W_{\mathbb R}^\perp}$, $\lambda^*|_{S_+\otimes W^*}$, $\lambda^*|_{S_-\otimes W}$, $c^*$, are those of Proposition \ref{prop:c11-oddcensus} (ii).\label{tab:c11-defectcounts}}

\begin{proposition}[The quadratic free directions along $\Sigma$ and $\Sigma^*$]
\label{prop:c11-freexi}
\begin{itemize}
\item (i) The quadratic local functionals of degree $-1$, of Weyl weight $\tfrac92$ in the fields and derivatives, linear in $\xi\in S_+\otimes W$, translation invariant and covariant under $\mathrm{Spin}(4)\times\mathrm{Spin}(6)\times\mathfrak g$, are spanned by $T_1$, $T_2$, $T_3$, $T_4$ and $\int\big((d^*A|_V)\xi,\lambda^*\big)$; their $Q$-closed combinations are spanned by $S^V_2(\xi)$ and $S^\perp_2(\xi)$, and the $Q$-exact ones are $0$.
\item (ii) The arity-one equation along $\xi$, the arity-one part along $\xi$ of $s(\hbar F_1)=\hbar\,\mathrm{Ob}$ (Theorem \ref{thm:c8-obstruction} (ii)), is
\begin{equation}
(Q+\{I,-\})\,\hbar F_1^{(1)}(\xi)+\big\{S^{(1)}(\xi),\hbar F_1^{(0)}\big\}=\hbar\,\mathrm{Ob}^{(1)}(\xi).
\label{eq:c11-arity1xi}
\end{equation}
At fixed $\hbar F_1^{(0)}$ two solutions differ, in the quadratic sector, by $\theta^\xi_V$ units times $S^V_2(\xi)$ plus $\theta^\xi_\perp$ units times $S^\perp_2(\xi)$, and the cubic part of \eqref{eq:c11-arity1xi} then contains $\theta^\xi_V\{I_3,S^V_2(\xi)\}+\theta^\xi_\perp\{I_3,S^\perp_2(\xi)\}$ units, $I_3$ the cubic part of $I$. For $\theta^\xi_V=\theta^\xi_\perp=:\theta^\xi$ the direction extends to $\theta^\xi$ units times $S^{(1)}(\xi)$, the arity-one value of $G_B$ (Definition \ref{def:c8-B}) and of the field redefinition by the R-rotation $S^{(1)}(R_\xi)$, $[R_\xi,\xi]=\xi$ (Theorem \ref{thm:c8-twist}). The cubic part fixes $\theta^\xi_V-\theta^\xi_\perp$ when the class of $\{I_3,S^V_2(\xi)\}$ modulo the images under $Q$ of the cubic local functionals of degree $-1$ is nonzero, as for the sector coefficients along $\Delta$ (Proposition \ref{prop:c11-arity1} (iii)).
\item (iii) The one-body kernel of a letter $X$ along $\xi$ is $-\{\hbar F_1^{(1)}(\xi),X\}$; the free directions add $-\theta^\xi_V$ units times the arrows of $S^V_2(\xi)$ and $-\theta^\xi_\perp$ units times those of $S^\perp_2(\xi)$, which are $\{S^V_2(\xi),A|_V\}=\Gamma(\xi,\lambda)|_V$ and $\{S^\perp_2(\xi),A|_{W_{\mathbb R}^\perp}\}=\Gamma(\xi,\lambda)|_{W_{\mathbb R}^\perp}$ on the gauge field and the scalars and $\tfrac12\rho(dA|_V)\xi$ and $\tfrac12\rho(dA|_{W_{\mathbb R}^\perp})\xi$ on $\lambda$ (Lemma \ref{lem:c11-lines} (i)).
\item (iv) Along $\Sigma^*$, with the position included at weight $-1$ and $\eta$ at weight $\tfrac12$, the $Q$-closed quadratic directions of degree $-1$ and Weyl weight $4$, linear in $\eta$, are spanned by $\sum_\mu\int x^\mu S^V_2(\rho(e_\mu)\eta)$ and $\sum_\mu\int x^\mu S^\perp_2(\rho(e_\mu)\eta)+\int s_0$, whose sum is the quadratic part of $S^{(1)}(\eta)$.
\end{itemize}
The same holds along $S_-\otimes W^*$ with the chiralities exchanged.
\end{proposition}

\begin{proof}
(i) The field contents of degree $-1$ and field weight $\tfrac92$ linear in a constant spinor are $(\lambda,A^*)$ without derivatives and $(A,\lambda^*)$ with one derivative, since $c$ and $c^*$ would need a further weight $\tfrac12$; the covariant pairings are $\Gamma(\xi,\lambda)$ with $A^*$ in its parts along $V$ and $W_{\mathbb R}^\perp$, and $\rho(\partial A)\xi$ with $\lambda^*$, $\partial A$ being $dA|_V$, $d^*A|_V$ or the mixed part $dA|_{W_{\mathbb R}^\perp}$, the volume element acting by a sign on each chirality. On functionals $Q$ substitutes $A\mapsto-dc$, $A^*\mapsto-\tfrac12d^*dA$, $\lambda^*\mapsto{\not d}\lambda$, $c^*\mapsto d^*(A^*|_V)$ and annihilates $c$ and $\lambda$ (Proposition \ref{prop:c1-Q}). $Q\int((d^*A|_V)\xi,\lambda^*)$ contains $-\int((d^*dc)\xi,\lambda^*)$, produced by no other structure, so its coefficient vanishes in a closed combination; $Q$ preserves the split into the parts along $V$ and along $W_{\mathbb R}^\perp$, and $S^V_2(\xi)$, $S^\perp_2(\xi)$ are closed, being the two parts of the quadratic part of $\{S,S^{(1)}(\xi)\}=0$. A $Q$-exact functional is $Q$ of a functional of degree $-2$ and field weight $\tfrac92$, and two antifields have weight at least $5$, while $c^*$ needs a further field of weight $\tfrac12$. Both evaluations of Remark \ref{rem:c11-oddmethod} give the same ranks on plane-wave normal forms. (ii) Theorem \ref{thm:c8-obstruction} (iii) restricted to arity one along $\xi$ at fixed arity-zero part. The cubic part of $s\big(\theta^\xi_VS^V_2+\theta^\xi_\perp S^\perp_2\big)$ is the bracket with $I_3$, and $\{I_3,S^V_2(\xi)+S^\perp_2(\xi)\}=-QS_3(\xi)$ with $S_3(\xi)$ the cubic part of $S^{(1)}(\xi)$, the cubic part of $\{S,S^{(1)}(\xi)\}=0$; $G_B^{(1)}(\xi)=S^{(1)}(B\xi)=S^{(1)}(\xi)$ for $\xi\in S_+\otimes W$. (iii) Theorem \ref{thm:c11-classes} (iii) and Lemma \ref{lem:c11-lines} (i). (iv) The same enumeration with $\int x^\mu$ times the structures of (i) at $\rho(e_\mu)\eta$, $\int s_0$ and $\int(\rho(A|_V)\eta,\lambda^*)$, and the exact ranks of $Q$ on plane-wave normal forms; $s_0$ enters the closed direction together with the part along $W_{\mathbb R}^\perp$.
\end{proof}

\begin{proposition}[Renormalisation freedom of $Q_1$, $S_1$ and $K_1$]
\label{prop:c11-freedom}
Under the changes of Proposition \ref{prop:c8-laws} the components change as follows; $\{S^{(1)}(x),-\}$ is the classical transformation along $x$.
\begin{itemize}
\item (i) A field redefinition $F_{\mathfrak{psl}(4|4)}\mapsto F_{\mathfrak{psl}(4|4)}+\hbar s\phi$ changes the component along $x$ by $-\{(s\phi)^{(1)}(x),-\}$, a derivation of the letters: it changes the one-body kernels and leaves the two-body kernels unchanged. For the field countings $\phi^{(0)}=\sum_fy_f\int(f,f^*)$ the change of $Q_1$ multiplies each arrow $f\to f'$ of the classical transformation by $y_{f'}-y_f$; for $\phi^{(0)}=t\,S^{(1)}(\Delta)$ it is a multiple of the classical transformation proportional to $\lambda_x$.
\item (ii) The coupling direction leaves $Q_1$, $S_1$, $K_1$ unchanged. The outer direction $G_B$ changes $Q_1$ along $S_+\otimes W$ by $-a$ and along $S_-\otimes W^*$ by $+a$ times the classical transformation, $S_1$ along $S_+\otimes W^*$ by $+a$ and along $S_-\otimes W$ by $-a$ times it, and leaves $K_1$, $D_1$, $P_1$, $M_1$, $R_1$ unchanged.
\item (iii) The free directions of $\hbar F_1^{(1)}(x)$ change the one-body kernels: along $\Sigma$ by Proposition \ref{prop:c11-freexi} (iii); along $\Delta$ and $P_b$ through the sector coefficients of Definition \ref{def:c11-free}, which enter $D_1$, $P_1$ and, through $P_1$, the relations of Proposition \ref{prop:c11-S} and Corollary \ref{cor:c11-K}; along $K_b$ through the $s$-closed directions of degree $-1$ quadratic in the Killing data. $\hbar F_1^{(0)}$ enters the one-body kernels through \eqref{eq:c11-arity1xi} and through the trees with $S^{(1)}(x)$ (Remark \ref{rem:c10-scheme1} (a), (b)).
\item (iv) A finite part $\psi'$ of the operators changes the component along $x$ by $[\{S^{(1)}(x),-\},\psi']$, on $O$ by $\{S^{(1)}(x),\psi'(O)\}-\psi'(\{S^{(1)}(x),O\})$. For $\psi'=\tfrac a2\{S^{(1)}(\Delta),-\}$ this is $\tfrac a2\lambda_x\{S^{(1)}(x),-\}$: $+a$, $-\tfrac a2$, $+\tfrac a2$ times the classical transformation along $K_b$, $\Sigma$, $\Sigma^*$ (Proposition \ref{prop:c9-scheme}), a change of the one-body kernels. A finite part $\psi'$ that vanishes on $O$ changes the component of $O$ along $x$ by $-\psi'(\{S^{(1)}(x),O\})$: the finite parts of the operators of the Weyl weight of $\{S^{(1)}(x),O\}$ enter the two-body kernels, as the finite part of Lemma \ref{lem:c11-cocycle} on the trace bilinear and the shift of the Konishi entry in Remark \ref{rem:c11-twist}.
\item (v) The cutoff shape enters the finite parts of the defect cycles (Proposition \ref{prop:c10-shape}, Lemma \ref{lem:c10-special}), of the special-edge cycles of $L_1^{(0)}$ and of the finite special vertices; since the counterterms are unique up to finite local terms [C07, Thm B], a change of the shape is a change of $\hbar F_1$ and of $\psi$, a combination of (i)--(iv).
\item (vi) Under the hypotheses of Theorem \ref{thm:c11-arity2}, in the translation-covariant scheme at $\theta=\theta_\lambda=\theta_V=-\tfrac13$, $P_1\cong0$ on the families of Definition \ref{def:c11-families} (Corollary \ref{cor:c11-D1fam}), and the condition $P_1=0$ of Definition \ref{def:c8-radial} holds there with $\psi'=0$: the components along $K_b$, $\Sigma$ and $\Sigma^*$ are those of the scheme. For equal coefficients $\theta_P\ne\theta$ along $P_b$, the finite part with $a=(\theta-\theta_P)$ units gives $P_1\cong0$ and changes the components along $K_b$, $\Sigma$, $\Sigma^*$ by $a$, $-\tfrac a2$, $\tfrac a2$ times the classical transformation (Corollary \ref{cor:c11-theta} (iv)).
\item (vii) Under (i), (iv) and the coupling direction the representation of $\mathfrak{psl}(4|4)$ on the cohomology of $\Lambda^{(0)}$ changes by an isomorphism, under $G_B$ by the automorphism $e^{-a\hbar B}$, and the spectrum of $D_1$ in a fixed classical Weyl weight is invariant (Theorem \ref{thm:c8-invariants}); the class $[Q^\xi_1]$ on the cohomology of $Q^\xi_0$, its kernel, image and rank are invariant under (i), (ii) and (iv) (Theorem \ref{thm:c8-twist}).
\end{itemize}
\end{proposition}

\begin{proof}
(i), (ii) Proposition \ref{prop:c8-laws} (a), (c); the bracket with a local functional is a derivation in its second argument, so it acts letter by letter; for $\phi^{(0)}=t\,S^{(1)}(\Delta)$, $\{S^{(1)}(x),S^{(1)}(\Delta)\}$ is $S^{(1)}$ of $\pm[x,\Delta]=\mp\lambda_xx$ by Lemma \ref{lem:c2-components}, with $S^{(2)}(x,\Delta)=0$ (Definition \ref{def:c2-S2}). (iii) Proposition \ref{prop:c11-freexi}, Definition \ref{def:c11-free}, Remark \ref{rem:c10-scheme1}. (iv) Proposition \ref{prop:c8-laws} (b) and Proposition \ref{prop:c9-scheme}, with $[\{S^{(1)}(x),-\},\psi'](O)=\{S^{(1)}(x),\psi'(O)\}-\psi'(\{S^{(1)}(x),O\})$. (v) Remark \ref{rem:c10-scheme1} (d); the finite counterterm and the finite parts of the operators absorb the change. (vi) Corollaries \ref{cor:c11-D1fam} and \ref{cor:c11-theta} (iv). (vii) Theorems \ref{thm:c8-invariants} and \ref{thm:c8-twist}.
\end{proof}

\begin{lemma}[The trace bilinear in a cocycle scheme]
\label{lem:c11-cocycle}
Let $O=\int\omega\,(\iota_{w_1}A,\iota_{w_2}A)$, $\xi\in S_+\otimes W$ or $\xi\in S_-\otimes W^*$, and let $\mathcal B_1,\dots,\mathcal B_4$ be the four functionals of Corollary \ref{cor:c7-Q},
\[
\mathcal B_1:=\{S^{(1)}(\xi),O\},\quad \mathcal B_2:=(w_1,w_2)\!\int\!\omega\,\big(\rho(A|_{W_{\mathbb R}^\perp})\xi,\lambda\big),\quad \mathcal B_3:=(w_1,w_2)\!\int\!\omega\,\big(\rho(A|_V)\xi,\lambda\big),\quad \mathcal B_4:=(w_1,w_2)\!\int\!\omega\,\big(c,(\xi,\lambda^*)\big),
\]
$\mathcal B_1$ being the first functional listed there (Lemma \ref{lem:c11-lines} (i)). Write the two-body part of $Q_1(O)$ as $a_1\mathcal B_1+a_2\mathcal B_2+a_3\mathcal B_3+a_4\mathcal B_4$.
\begin{itemize}
\item (i) If $L_1^{(0)}=0$ on $O$ and on $\mathcal B_1$, then $a_3=a_4=0$, and the one-body terms of $Q_1(O)$ are multiples of $\mathcal B_1$.
\item (ii) If $L_1^{(0)}(O)=2\kappa_O(w_1,w_2)\int\omega\,(dc,A|_V)$ and the part of $L_1^{(0)}(\mathcal B_1)$ with the two letters $c$, $\lambda$ is $-\kappa_{\mathcal B}\,s\mathcal B_3$, the finite part $\psi(O)=\kappa_O(w_1,w_2)\int\omega\,(A|_V,A|_V)$, $\psi(\mathcal B_1)=\kappa_{\mathcal B}\mathcal B_3$ removes both and changes $Q_1(O)$ by $(2\kappa_O-\kappa_{\mathcal B})\mathcal B_3$. At $\psi=0$ the part with the letters $c$, $\lambda$ of (R1) of Proposition \ref{prop:c11-relations} is the identity $a_3+2\kappa_O-\kappa_{\mathcal B}=0$.
\end{itemize}
\end{lemma}

\begin{proof}
(i) By (R1), $s\,Q_1(O)=0$. $s\mathcal B_1=0$, since $[s,\{S^{(1)}(\xi),-\}]=\{\{S,S^{(1)}(\xi)\},-\}=0$ at tree level and $sO=0$; $s\mathcal B_2=0$, since $Q$ annihilates the coordinate functionals of $A|_{W_{\mathbb R}^\perp}$ and $\lambda$ and the cubic parts of $s$ on them are the gauge transformation with parameter $c$, which leaves the invariant pairing fixed. In the sector with the two letters $c$, $\lambda$, $s\mathcal B_3=-(w_1,w_2)\int\omega\,(\rho(dc)\xi,\lambda)$ and $s\mathcal B_4=(w_1,w_2)\int\omega\,(c,(\xi,{\not d}\lambda))$, whose densities carry the derivative on different letters and are linearly independent for arbitrary $\omega$; hence $a_3=a_4=0$. The one-body terms are $-\{\hbar F_1^{(1)}(\xi),O\}$; a term of $\hbar F_1^{(1)}(\xi)$ acting on the letters $A|_{W_{\mathbb R}^\perp}$ of $O$ has the leg $A^*|_{W_{\mathbb R}^\perp}$, and by Proposition \ref{prop:c11-freexi} (i) it is a multiple of $T_2$, which gives a multiple of $\mathcal B_1$. (ii) $Q\int\omega\,(A|_V,A|_V)=-2\int\omega\,(dc,A|_V)$ (Remark \ref{rem:c5-cocycle}), and the cubic part of $s\int\omega\,(A|_V,A|_V)$ vanishes by the invariance of $(-,-)$; $L_1$ changes by $[s,\psi]$ (Lemma \ref{lem:c10-psi}), the component along $\xi$ by $[\{S^{(1)}(\xi),-\},\psi]$ (Proposition \ref{prop:c8-laws} (b)), and $\{S^{(1)}(\xi),\int\omega\,(A|_V,A|_V)\}=2\int\omega\,(\Gamma(\xi,\lambda)|_V,A|_V)$ with $(\Gamma(\xi,\lambda)|_V,a)=(\rho(a)\xi,\lambda)$ [ESW, (28)]. In the shifted scheme (i) gives $a_3+2\kappa_O-\kappa_{\mathcal B}=0$; at $\psi=0$ the part with the letters $c$, $\lambda$ of (R1), $s\,Q_1(O)+L_1^{(0)}(\mathcal B_1)+\{S^{(1)}(\xi),L_1^{(0)}(O)\}=0$, is the same equation, since there $s\,Q_1(O)=a_3\,s\mathcal B_3$, $L_1^{(0)}(\mathcal B_1)=-\kappa_{\mathcal B}\,s\mathcal B_3$ and $\{S^{(1)}(\xi),L_1^{(0)}(O)\}=2\kappa_O\,s\mathcal B_3$.
\end{proof}

\begin{theorem}[$Q_1$ on the scalar bilinears]
\label{thm:c11-bilinears}
Let $O=\int\omega\,(\iota_{w_1}A,\iota_{w_2}A)$, $\xi\in S_+\otimes W$ or $\xi\in S_-\otimes W^*$, and $\mathcal B_1,\dots,\mathcal B_4$ as in Lemma \ref{lem:c11-cocycle}. At $\psi=0$:
\begin{itemize}
\item (i) the two-body part of $Q_1(O)$ is
\begin{equation}
-2L\,\mathcal B_1+4L\,\mathcal B_2+4L\,\mathcal B_3+0\cdot\mathcal B_4\ \text{units},\qquad L=\log\tfrac32,
\label{eq:c11-Q1O}
\end{equation}
without poles and the same for both chiralities; its trace part is $4L$ units times $(w_1,w_2)\int\omega\,(\rho(A)\xi,\lambda)$, with the whole connection $A=A|_V+A|_{W_{\mathbb R}^\perp}$;
\item (ii) on the traceless part, $(w_1,w_2)=0$,
\[
Q_1(O)=c_Q\,\{S^{(1)}(\xi),O\},\qquad c_Q=-2\log\tfrac32-t_\perp\ \text{units},
\]
with $t_\perp$ units the coefficient of $T_2$ in $\hbar F_1^{(1)}(\xi)$, which \eqref{eq:c11-arity1xi} fixes up to the free direction $\theta^\xi_\perp$ of Proposition \ref{prop:c11-freexi}; $c_Q$ moves by $-\theta^\xi_\perp$ under that direction, by $\mp a$ under $G_B$ along $S_+\otimes W$ and $S_-\otimes W^*$, by $-\tfrac a2$ under $\psi'=\tfrac a2\{S^{(1)}(\Delta),-\}$, and under the field countings as in Proposition \ref{prop:c11-freedom} (i);
\item (iii) $L_1^{(0)}(O)=-(2+4L)$ units times $(w_1,w_2)\int\omega\,(dc,A|_V)$, and the part of $L_1^{(0)}(\mathcal B_1)$ with the letters $c$, $\lambda$ is $2$ units times $s\mathcal B_3$; Lemma \ref{lem:c11-cocycle} (ii) applies with $\kappa_O=-(1+2L)$ units and $\kappa_{\mathcal B}=-2$ units, and in the cocycle scheme
\[
\psi(O)=-(1+2L)\,\text{units}\cdot(w_1,w_2)\int\omega\,(A|_V,A|_V),\qquad \psi(\mathcal B_1)=-2\,\text{units}\cdot\mathcal B_3,
\]
the two-body part of $Q_1(O)$ is $-2L\,\mathcal B_1+4L\,\mathcal B_2$ units, $a_3=4L+2\kappa_O-\kappa_{\mathcal B}=0$;
\item (iv) at $\psi=0$, for $\xi\in S_+\otimes W$, the part with the letters $c$, $\lambda|_{S_-\otimes W^*}$ of (R1) is the identity
\[
s\,Q_1(O)+L_1^{(0)}\big(\{S^{(1)}(\xi),O\}\big)+\big\{S^{(1)}(\xi),L_1^{(0)}(O)\big\}=\big(-4L-2+(2+4L)\big)\,\text{units}\cdot(w_1,w_2)\int\omega\,\big(\rho(dc)\xi,\lambda\big)=0 .
\]
\end{itemize}
\end{theorem}

\begin{proof}
(i) By Theorem \ref{thm:c11-classes} (iii) the two-body part is $\Upsilon_\xi$ on the single pair $(A|_{W_{\mathbb R}^\perp},A|_{W_{\mathbb R}^\perp})$ of $O$, whose leg sets are those of Table \ref{tab:c11-defectpairs}: $\lambda|_{S_-\otimes W^*}$, $A|_V\,\lambda|_{S_-\otimes W^*}$, $A|_{W_{\mathbb R}^\perp}\lambda|_{S_+\otimes W}$ and $c\,\lambda^*|_{S_+\otimes W^*}$ for $\xi\in S_+\otimes W$. The single output $Y$ has the colour factor $(1\otimes\mathrm{ad}_Y)(-,-)^{-1}$, which the colour pairing of $O$ contracts to $\operatorname{tr}_{\mathfrak g}\mathrm{ad}_Y=0$. The other three leg sets have $D=0$, so they carry no poles and no finite special vertex (Proposition \ref{prop:c11-oddbound}), and they give $\mathcal B_1$, $\mathcal B_2$ (outputs $A|_{W_{\mathbb R}^\perp}\lambda$), $\mathcal B_3$ (outputs $A|_V\lambda$) and $\mathcal B_4$ (outputs $c\,\lambda^*$). Their chains have three edges and, at $D=0$, constant integrands, so each chain gives a rational multiple of $\tfrac12\log\tfrac32$ (Theorem \ref{thm:c11-oddsimplex} (iv)); the exact evaluation of the chains in all placements, in both ways of Remark \ref{rem:c11-oddmethod}, gives \eqref{eq:c11-Q1O}, with the same coefficient for $\mathcal B_2$ and $\mathcal B_3$. (ii) By Corollary \ref{cor:c7-Q} the functionals at $(w_1,w_2)=0$ are the multiples of $\mathcal B_1$; the one-body term is $-t_\perp\mathcal B_1$ units (proof of Lemma \ref{lem:c11-cocycle} (i)); the changes are Propositions \ref{prop:c11-freexi} and \ref{prop:c11-freedom}. (iii) $L_1^{(0)}$ is given by the arity-zero two-body kernels, the cycles with one special edge of the proof of Theorem \ref{thm:c10-brst}, of the pairs $(A|_{W_{\mathbb R}^\perp},A|_{W_{\mathbb R}^\perp})$ of $O$ and $(A|_{W_{\mathbb R}^\perp},\lambda)$ of $\mathcal B_1$, evaluated in the same way; then Lemma \ref{lem:c11-cocycle} (ii). (iv) With $s\mathcal B_3=-(w_1,w_2)\int\omega\,(\rho(dc)\xi,\lambda)$ in this sector, the three terms are $4L\,s\mathcal B_3$, $2\,s\mathcal B_3$ and $2\kappa_O\,s\mathcal B_3=-(2+4L)\,s\mathcal B_3$ units (proof of Lemma \ref{lem:c11-cocycle} (ii)).
\end{proof}

\begin{proposition}[$S_1$ and $K_1$ on the scalar bilinears]
\label{prop:c11-SKbil}
Let $O=\int\omega\,(\iota_{w_1}A,\iota_{w_2}A)$ and $\psi=0$.
\begin{itemize}
\item (i) $\Upsilon^0_\eta(A|_{W_{\mathbb R}^\perp},A|_{W_{\mathbb R}^\perp})$ has the single leg set $\lambda$, which vanishes on the colour singlet, and $s_0$ has zero bracket with a letter $A|_{W_{\mathbb R}^\perp}$. Hence the two-body part of $S_1(O)$ at the point $x$ of the density is the two-body part of $Q_1(O)$ along the Killing spinor $\rho(x)\eta$, given by \eqref{eq:c11-Q1O} along $S_-\otimes W^*$ for $\eta\in S_+\otimes W^*$ and along $S_+\otimes W$ for $\eta\in S_-\otimes W$; the term linear in ${\not d}(\rho(x)\eta)=4\eta$ is absent (Corollary \ref{cor:c7-S}).
\item (ii) $\Upsilon^0_{K_b}(A|_{W_{\mathbb R}^\perp},A|_{W_{\mathbb R}^\perp})$ has the leg sets $\emptyset$, of superficial degree $1$, and $A|_V$, both zero on the colour singlet. In a scheme in which $\hbar F_1^{(1)}(K_b)$ decomposes as $S^{(1)}(K_b)$ in Theorem \ref{thm:c11-oddsimplex} (iii), on the part with two legs $A|_{W_{\mathbb R}^\perp}$,
\[
K_1(O)=v(x)^\mu\,P_1(O)\big|_{P_{e_\mu}}-(b,x)\,D^0_1(O),
\]
with $D^0_1$ the translation-invariant part of $D_1$ (Definition \ref{def:c11-canonical}); the rotation part and the intrinsic part vanish. With Theorems \ref{thm:c9-D1} and \ref{thm:c9-P1} at $\theta_P=\theta$,
\[
K_1(O)=-\big(\tfrac13+\theta\big)\,\text{units}\cdot\{S^{(1)}(K_b),O\}+\begin{cases}0&\text{(traceless part)},\\ 48\,\text{units}\cdot(b,x)\,O&\text{(trace)},\end{cases}\qquad \{S^{(1)}(K_b),O\}=v^\mu\partial_\mu O-4(b,x)\,O,
\]
so that at $\theta=-\tfrac13$, $K_1(O)=-(b,x)D_1(O)$: $0$ on the traceless bilinears and $48$ units times $(b,x)\,O$ on the trace, the relation of Proposition \ref{prop:c8-primary} with $\kappa=-24$ units.
\end{itemize}
\end{proposition}

\begin{proof}
(i) Table \ref{tab:c11-defectcounts} and the evaluation of the leg set: a single output $Y$ gives $\operatorname{tr}_{\mathfrak g}\mathrm{ad}_Y=0$ on the colour singlet; $s_0$ has the legs $A|_{W_{\mathbb R}^\perp}$, $\lambda^*$, and its bracket with a letter needs the leg $A^*|_{W_{\mathbb R}^\perp}$. Theorem \ref{thm:c11-oddsimplex} (ii) gives the position part. (ii) The leg set $A|_V$ has a single output; the empty leg set is homogeneous in $\varepsilon$ of degree $-\tfrac12$ (Lemma \ref{lem:c10-self}). Theorem \ref{thm:c11-oddsimplex} (iii) for the two-body parts, with the one-body term decomposed in the same way; by Corollary \ref{cor:c7-K} the term in the second derivatives of $v$ and the term in the antisymmetric part of $dv$ vanish on a scalar bilinear. By \eqref{eq:c9-P1}, $P_1(O)|_{P_{e_\mu}}=-(\tfrac13+\theta_P)$ units times $\partial_\mu O$; by \eqref{eq:c9-D1}, with $\{S^{(1)}(\Delta),O\}=2x^\mu\partial_\mu O+4O$, $D^0_1(O)$ at the point of the density is $-4(\tfrac13+\theta)$ units times $O$, plus $-48$ units times $O$ on the trace; and $\{S^{(1)}(K_b),O\}=v^\mu\partial_\mu O+\tfrac24(\operatorname{div}v)O$ (Definition \ref{def:c7-vdot}). Proposition \ref{prop:c8-primary}: $K_1(O)=-2(b,x)\kappa\,O$ for $D_1(O)=2\kappa\,O$.
\end{proof}

\begin{corollary}[The dilatation and the translations by the defect cycles]
\label{cor:c11-DPdefect}
On every local functional $O$ with $sO=0$, with the Koszul signs of \eqref{eq:c10-decomp},
\[
\sum_{a<b}\pm\Upsilon_\Delta(X_a,X_b)=\sum_{a<b}\pm\big(4\,\mathrm{lc}(X_aX_b)+\text{composite pair kernel along }\Delta\big),\qquad \sum_{a<b}\pm\Upsilon_{P_b}(X_a,X_b)=\sum_{a<b}\pm\big(\text{composite pair kernel along }P_b\big),
\]
and the one-body terms of the two forms agree. On the part with two legs $A|_{W_{\mathbb R}^\perp}$ of the scalar bilinears, at the point of the density, $\Upsilon^0_\Delta$ is $\tfrac{40}3$ units times $O$ on the traceless part and $-\tfrac{104}3$ units times $O$ on the trace, and $\Upsilon_{P_b}$ is $-\tfrac13$ units times $\partial_bO$; with the position part $2x^\mu\Upsilon_{P_{e_\mu}}$ this is $\tfrac{44}3O-\tfrac13\{S^{(1)}(\Delta),O\}$ and $-\tfrac{100}3O-\tfrac13\{S^{(1)}(\Delta),O\}$ units, the first row of the table of Theorem \ref{thm:c9-D1}, and the value $-\tfrac13$ of Theorem \ref{thm:c9-P1}.
\end{corollary}

\begin{proof}
Both forms compute the unique local representative (Theorem \ref{thm:c8-L} (i)): for $\Delta$ this is Proposition \ref{prop:c11-routes}, whose route of the special edges is \eqref{eq:c11-decompx}, and for $P_b$ the same argument applies with Corollary \ref{cor:c11-D1local} and the one-body term $-\{\hbar F_1^{(1)}(P_b),-\}$ in both forms. On the bilinears, the defect cycles with $S^{(1)}(\Delta)$ and $S^{(1)}(P_b)$ in all their placements, evaluated by Theorem \ref{thm:c11-oddsimplex}, give the stated values, with $\{S^{(1)}(\Delta),O\}=2x^\mu\partial_\mu O+4O$ at the point of the density: $\tfrac{40}3+\tfrac43=\tfrac{44}3$ and $-\tfrac{104}3+\tfrac43=-\tfrac{100}3$.
\end{proof}

\begin{proposition}[The special supercharges and the supercharges]
\label{prop:c11-S}
Let $\eta\in\Sigma^*$ and $b\in V$, so that $\rho(b)\eta\in\Sigma$.
\begin{itemize}
\item (i) Graph by graph, the position part along $b$ of the two-body kernel of $(X_a,X_b)$ along $\eta$ is the two-body kernel of $(X_a,X_b)$ along $\rho(b)\eta$; its translation-invariant part is computed with $S^{(1)}(\eta)$ at the spinor $\rho(x-x_O)\eta$, $x_O$ the point of the density, together with the term $-\int(\rho(A|_{W_{\mathbb R}^\perp})\eta,\lambda^*)$.
\item (ii) Under the hypotheses of Proposition \ref{prop:c7-select} and Corollary \ref{cor:c8-homotopy}, on cohomology, the position part along $b$ of $S_1$ along $\eta$ is $Q_1$ along $\rho(b)\eta$ minus the commutator of $\{S^{(1)}(\eta),-\}$ with $P_1$.
\item (iii) On a family on which $P_1\cong0$, in particular, under the hypotheses of Theorem \ref{thm:c11-arity2}, on the families of Definition \ref{def:c11-families} in the translation-covariant scheme at $\theta=\theta_\lambda=\theta_V=-\tfrac13$, the position part of $S_1$ along $\eta$ is $Q_1$ along $\rho(b)\eta$ on cohomology, and $S_1$ along $\eta$ is fixed by $Q_1$ along $\Sigma$ up to the term of Proposition \ref{prop:c7-select} linear in ${\not d}(\rho(x)\eta)=4\eta$. On the scalar bilinears this term is absent (Corollary \ref{cor:c7-S}), and $S_1$ along $\eta$ is, on cohomology, the combination of the functionals of Corollary \ref{cor:c7-S} with the coefficients of $Q_1$ along $\Sigma$ in Corollary \ref{cor:c7-Q}, given by Theorem \ref{thm:c11-bilinears} and Proposition \ref{prop:c11-SKbil} (i).
\item (iv) In a reflection-compatible scheme, on the family formed by a traceless scalar bilinear $O$ at the origin and its descendants $\{S^{(1)}(\xi),O\}$, on which $Q_1=c_Q\{S^{(1)}(\xi),-\}$ along the unit generator $Q^i_\alpha$ (Theorem \ref{thm:c11-bilinears} (ii)) and $S_1=c_S\{S^{(1)}(\eta),-\}$ along $S^\alpha{}_i$, Corollary \ref{cor:c8-gram} reads
\[
c_S=\overline{c_Q}+g_{5/2}-g_2,
\]
with $g_w$ the eigenvalue of $\mathrm{Gram}_{\rm tree}^{-1}\mathrm{Gram}_{\rm 1}$ in Weyl weight $w$ on the family; the outer direction changes $c_Q$ by $-a$ and $c_S$ by $+a$ and preserves the relation for $\bar a=-a$.
\end{itemize}
\end{proposition}

\begin{proof}
(i) By \eqref{eq:c11-Seta}, $S^{(1)}(\eta)=\int x^\mu s_\mu+\int s_0$, where $\int s_\mu$ is $S^{(1)}$ along the constant spinor $\rho(e_\mu)\eta\in\Sigma$. Insert $x^\mu=x_O^\mu+(x-x_O)^\mu$ and argue as in the proof of Proposition \ref{prop:c11-position} (i): on plane waves the factor $x^\mu$ is $-i\partial/\partial r^\mu$ of the momentum injected at the vertex, and its action on the momentum delta function is $-i$ times the graph with $\int s_\mu$; the same holds for the special edges, the vertices $S^{(1)}(\eta)$ on the cycle and the counterterm vertices, which are the defect insertions of Proposition \ref{prop:c10-arity1} (Theorem \ref{thm:c11-oddsimplex} (ii)). (ii) Corollary \ref{cor:c8-homotopy} (ii) at $(\eta,P_b)$: by Proposition \ref{prop:c2-special}, $[P_b,\eta]=\rho(b)\eta$, so $[\eta,P_b]=-\rho(b)\eta$, and $S^{(2)}(\eta,P_b)=0$ (Definition \ref{def:c2-S2}). At order $\hbar$ the identity states that the commutator of $S_1$ along $\eta$ with $\{S^{(1)}(P_b),-\}$, plus the commutator of $\{S^{(1)}(\eta),-\}$ with $P_1$, minus $Q_1$ along $\rho(b)\eta$, is a graded commutator with $\Lambda^{(0)}$, which induces $0$ on cohomology; at order $\hbar^0$ the same identity is the statement of (i) for the classical transformation, which fixes the signs. By Proposition \ref{prop:c10-bound1}, $S_1$ along $\eta$ is at most linear in the position, and Lemma \ref{lem:c11-canonical} identifies the first commutator with its position part. (iii) By Proposition \ref{prop:c7-select} the representative of $S_1$ along $\eta$ at a point $x_0$ is a term linear in the value $\rho(x_0)\eta$ plus a term linear in $4\eta$; the first has the position part given by its value at $\rho(b)\eta$, and as $\eta$ and $b$ vary, $\rho(b)\eta$ spans $\Sigma$, so by (ii) it is the representative of $Q_1$ evaluated at $\rho(x_0)\eta$ on cohomology. (iv) Write $X_0+\hbar X_1$ for the matrix of $-\Lambda^{(1)}(Q^i_\alpha)$ and $Y_0+\hbar Y_1$ for that of $-\Lambda^{(1)}(-S^\alpha{}_i)$ (Theorem \ref{thm:c8-dag} (i)): $\hbar X_1=c_QX_0$ and $\hbar Y_1=c_SY_0$, and $Y_0$ maps Weyl weight $\tfrac52$ to Weyl weight $2$. Corollary \ref{cor:c8-gram} gives $\hbar Y_1=\overline{c_Q}\,X_0^{\rm adj}+\hbar[Y_0,\mathrm{Gram}_{\rm tree}^{-1}\mathrm{Gram}_{\rm 1}]=\big(\overline{c_Q}+g_{5/2}-g_2\big)Y_0$ on the descendants, with $Y_0=X_0^{\rm adj}$. The outer direction: Proposition \ref{prop:c11-freedom} (ii) and Corollary \ref{cor:c8-gram}.
\end{proof}

\begin{corollary}[$K_1$ in the translation-covariant scheme]
\label{cor:c11-K}
Under the hypotheses of Theorems \ref{thm:c8-adjoint} and \ref{thm:c11-arity2}, on a family of Definition \ref{def:c11-families} at the origin of one length, in the translation-covariant scheme at $\theta=\theta_\lambda=\theta_V=-\tfrac13$, $K_1$ along $K_{\bar b}$ equals $\hbar\big[\{S^{(1)}(K_{\bar b}),-\},\mathrm{Gram}_{\rm tree}^{-1}\mathrm{Gram}_{\rm 1}\big]$, the statement $K=P^\dagger$ of Corollary \ref{cor:c8-gram}. For equal coefficients $\theta_P$ along $P_b$, $K_1$ along $K_{\bar b}$ is $(\theta-\theta_P)$ units times $\{S^{(1)}(K_{\bar b}),-\}$ plus the same commutator, with the radial Gram matrix of Definition \ref{def:c8-inner} in that scheme. On a scalar bilinear at the origin both sides vanish: the defect cycles give $K_1(O)=0$ there (Proposition \ref{prop:c11-SKbil} (ii)), and the commutator vanishes on a primary at the origin on which $\mathrm{Gram}_{\rm tree}^{-1}\mathrm{Gram}_{\rm 1}$ is diagonal. On its descendants $\partial_\mu O$ the relation equates the intrinsic kernels $\Upsilon^0_{K_b}$ of the descendant pairs with the difference of the two-point normalisations in Weyl weights $w_O+1$ and $w_O$ times the classical transformation.
\end{corollary}

\begin{proof}
Corollary \ref{cor:c11-D1fam} gives $P_1\cong0$, resp. $P_1\cong(\theta-\theta_P)$ units times $\{S^{(1)}(P_b),-\}$ with a real coefficient; Corollary \ref{cor:c8-gram} gives the two relations. At the origin $v(0)=0$ and $(b,0)=0$, so the formula of Proposition \ref{prop:c11-SKbil} (ii) vanishes; $\{S^{(1)}(K_{\bar b}),-\}$ lowers the Weyl weight by one and annihilates a primary at the origin.
\end{proof}

\begin{proposition}[The relations at order $\hbar$]
\label{prop:c11-relations}
Write $\Lambda^{(1)}(x)=\Lambda^{(1)}_0(x)+\hbar\Lambda^{(1)}_1(x)$ with $\Lambda^{(1)}_0(x)=\pm\{S^{(1)}(x),-\}$ (Theorem \ref{thm:c8-L} (iv)), so that $\hbar\Lambda^{(1)}_1(x)O$ is the component of $O$ along $x$, and $\Lambda^{(0)}=s+\hbar L_1^{(0)}$. Let $\mathfrak F$ be a family of colour-singlet local functionals of degree $0$, stable under $\Lambda$ and under the classical transformations, with $L_1^{(0)}=0$ on $\mathfrak F$ for a finite part $\psi$ as in Definition \ref{def:c10-hyp} (ii), and let $O\in\mathfrak F$ with $sO=0$.
\begin{itemize}
\item (R1) $s\big(\hbar\Lambda^{(1)}_1(x)O\big)=0$ for every generator $x$; along $\xi$, $s\,Q_1(O)=0$, an identity of local functionals.
\item (R2) For $\xi,\xi'\in\Sigma$ with $b=\Gamma(\xi,\xi')|_V$, $[\Lambda^{(1)}_1(\xi),\Lambda^{(1)}_0(\xi')]O+[\Lambda^{(1)}_0(\xi),\Lambda^{(1)}_1(\xi')]O+\Lambda^{(1)}_1(P_b)O=\Lambda^{(0)}(\chi)$, with $\chi$ local of degree $-1$, built from the arity-two component of $O$ along $(\xi,\xi')$; for $\xi,\xi'\in S_+\otimes W$, $b=0$.
\item (R3) For $\xi\in\Sigma$ and $\eta\in\Sigma^*$ with $[\xi,\eta]=\tfrac12(\xi,\eta)\Delta+M_{\xi\eta}+R_{\xi\eta}$ (Proposition \ref{prop:c2-special}), $M_{\xi\eta}\in\mathfrak{so}(V)$, $R_{\xi\eta}\in\mathfrak{so}(W_{\mathbb R}^\perp)$,
\[
[\Lambda^{(1)}_1(\xi),\Lambda^{(1)}_0(\eta)]O+[\Lambda^{(1)}_0(\xi),\Lambda^{(1)}_1(\eta)]O+\tfrac12(\xi,\eta)\,D_1(O)+M_1(O)+R_1(O)=\Lambda^{(0)}(\chi'),
\]
the components $M_1$, $R_1$ taken along $M_{\xi\eta}$, $R_{\xi\eta}$.
\item (R4) For a rank-one $\xi$, $Q^\xi_0Q^\xi_1+Q^\xi_1Q^\xi_0=0$ on every local functional (Theorem \ref{thm:c8-twist}).
\item (R5) With $[K_b,P_{b'}]=(b,b')\Delta+M_{bb'}$ (Remark \ref{rem:c7-KDelta}), $[\Lambda^{(1)}_1(K_b),\Lambda^{(1)}_0(P_{b'})]+[\Lambda^{(1)}_0(K_b),\Lambda^{(1)}_1(P_{b'})]+(b,b')\hbar\Lambda^{(1)}_1(\Delta)+\hbar\Lambda^{(1)}_1(M_{bb'})\cong0$ on cohomology; with Theorem \ref{thm:c11-oddsimplex} (iii) it states that the parts of $K_1$ multiplying the position are those of \eqref{eq:c11-Kdec}.
\item (R6) For $\eta_1,\eta_2\in\Sigma^*$ with $[\eta_1,\eta_2]=-K_b$, $b=\Gamma(\eta_1,\eta_2)|_V$, at a finite $T$, $K_1$ along $K_b$ is, on cohomology, $[\Lambda^{(1)}_1(\eta_1),\Lambda^{(1)}_0(\eta_2)]+[\Lambda^{(1)}_0(\eta_1),\Lambda^{(1)}_1(\eta_2)]$.
\end{itemize}
In (R2), (R3), (R5) and (R6) the right sides vanish on cohomology. For $O$ without letters $\lambda$ and $c$, such as the scalar bilinears and the families of scalars, the parts of order $\hbar^0$ of $\chi$, $\chi'$ vanish, the right sides are $s$-exact, and their parts with two field letters and no antifield vanish on free on-shell plane waves, $|p|^2=0$ on the scalars and the gauge field, $\rho(p)\lambda=0$ and $(p,a)=0$ on the polarisations $a$ of $A|_V$. On the scalar bilinears at $\theta=-\tfrac13$, (R3) contains $D_1(O)=0$ on the traceless part and $-48$ units times $O$ on the trace (Corollary \ref{cor:c9-ident}), the Konishi relation between $Q_1$, $S_1$ and the anomalous dimension, and (R2) with $\xi\in S_+\otimes W$, $\xi'\in S_-\otimes W^*$ contains $P_1(O)=-(\tfrac13+\theta_P)$ units times $\partial_bO$ (Theorem \ref{thm:c9-P1}).
\end{proposition}

\begin{proof}
Corollary \ref{cor:c8-homotopy} (i) at order $\hbar$ reads $[s,\Lambda^{(1)}_1(x)]+[L_1^{(0)},\Lambda^{(1)}_0(x)]=0$; on $O$ with $sO=0$, $L_1^{(0)}O=0$ and $L_1^{(0)}\{S^{(1)}(x),O\}=0$ (stability of $\mathfrak F$) this is (R1). Corollary \ref{cor:c8-homotopy} (ii) at order $\hbar$ for the pairs $(\xi,\xi')$, $(\xi,\eta)$, $(K_b,P_{b'})$, $(\eta_1,\eta_2)$, with the brackets of Proposition \ref{prop:c2-special} and Remark \ref{rem:c7-KDelta}, has the right side $[\Lambda^{(0)},H]$ with $H$ built from $\Lambda^{(2)}$, which on $O$ with $\Lambda^{(0)}O=0$ is $\Lambda^{(0)}(HO)$; the part of order $\hbar^0$ of $H$ is $\pm\{S^{(2)}(x_1,x_2),-\}$, and $S^{(2)}$ has the legs $(\lambda^*,\lambda^*)$ and $(A,c^*)$ (Definition \ref{def:c2-S2}), whose bracket with a functional of the fields $A$ alone vanishes. (R4) is Theorem \ref{thm:c8-twist}. For the vanishing on shell: the part of $s\chi$ with two field letters is $Q\chi_2+\{I_3,\chi_1\}$, with $\chi_k$ the part of $\chi$ with $k$ letters; $\chi_1$ is a single adjoint letter contracted with a constant tensor, zero for a colour singlet; $\chi_2$ of degree $-1$ has one antifield and one field, or $c^*$ and $c$, and $Q$ substitutes $A^*\mapsto-\tfrac12d^*dA$, $\lambda^*\mapsto{\not d}\lambda$, $c^*\mapsto d^*(A^*|_V)$, $A\mapsto-dc$ (Proposition \ref{prop:c1-Q}), so the part of $Q\chi_2$ without antifields and ghosts is a free field equation times a field. The values on the bilinears are Corollary \ref{cor:c9-ident} and Theorem \ref{thm:c9-P1}, with $R_1=0$ on them (Remark \ref{rem:c7-internal}).
\end{proof}

\begin{proposition}[The chiral-primary bilinears along a rank-one supercharge]
\label{prop:c11-plane}
Let $\xi\in S_+\otimes W$ be of rank one, $\xi_\alpha{}^i=u_\alpha z^i$, and let $w_1,w_2\in W_{\mathbb R}^\perp\otimes\mathbb C$ satisfy $\rho(w_1)\xi=\rho(w_2)\xi=0$. Then $(w_1,w_2)=0$, and under the hypotheses of Proposition \ref{prop:c7-select} the component along $\xi$ of $O=\int\omega\,(\iota_{w_1}A,\iota_{w_2}A)$ vanishes at every order at which they hold: its classical part $\{S^{(1)}(\xi),O\}$ and $Q_1(O)$ are $0$, in every scheme in which the counterterms are equivariant for translations, $\mathfrak{so}(V)$, $\mathfrak{so}(W_{\mathbb R}^\perp)$ and $\mathfrak g$.
\end{proposition}

\begin{proof}
$\rho(w_1)\rho(w_2)+\rho(w_2)\rho(w_1)=2(w_1,w_2)$ applied to $\xi\ne0$ gives $(w_1,w_2)=0$. By Corollary \ref{cor:c7-Q}, for $(w_1,w_2)=0$ the component is a multiple of $\int\omega\,\big[(\iota_{w_1}A,(\rho(w_2)\xi,\lambda))+(\iota_{w_2}A,(\rho(w_1)\xi,\lambda))\big]$, which vanishes; the classical component is the same functional (proof of Corollary \ref{cor:c7-Q}). In the evaluation of Theorem \ref{thm:c11-bilinears} the two-body part on these bilinears, $-2L\,\mathcal B_1+4L(\mathcal B_2+\mathcal B_3)$ units, vanishes term by term, with $\mathcal B_1=\{S^{(1)}(\xi),O\}=0$ and $(w_1,w_2)=0$, and the one-body term is a multiple of $\mathcal B_1$.
\end{proof}

\begin{remark}[The pair entries of the holomorphic twist]
\label{rem:c11-twist}
The twisted one-loop differential is $\hbar Q^\xi_1(O)=\hbar L_1(O)^{(0)}+Q_1(O)-\tfrac12\hbar L_1(O)^{(2)}(\xi,\xi)-\cdots$ (Definition \ref{def:c8-twisted}), so its pair rule [BK, (3.7)--(3.8)] corresponds to the two-body kernels of $Q_1$ along $\xi$ and of the arity-two component along $(\xi,\xi)$. By Proposition \ref{prop:c7-select}, applied once for each argument, a component of arity $k$ along $\xi$ of an operator of Weyl weight $w$ and degree $d$ has representatives of Weyl weight $w+\tfrac k2$ and degree $d+1-k$. [BK, (3.51)] relates the holomorphic letters to the fields up to field redefinitions; with $\xi_\alpha{}^i=u_\alpha z^i$, $u'\in S_+$ with $u'^\alpha u_\alpha=1$, and the representatives $\gamma\leftrightarrow A_{ij}z^j$ (Weyl weight $1$, degree $0$), $\beta\leftrightarrow u'^\alpha\lambda_\alpha{}^i$ modulo $z^i$ $(\tfrac32,0)$, $\partial_{\dot\alpha}c\leftrightarrow\bar\lambda_{\dot\alpha i}z^i$ $(\tfrac32,0)$, $c\leftrightarrow c$ $(0,1)$, $b\leftrightarrow u'^\alpha F_{mn}(\sigma^m\bar\sigma^n)_\alpha{}^\beta u'_\beta$ $(2,0)$ or a component of $\lambda^*$ $(\tfrac52,-1)$, and $\partial_{\dot\alpha}\leftrightarrow u'^\alpha\sigma^m_{\alpha\dot\alpha}D_m$ $(1,0)$, the pair entries of [BK, (3.64)--(3.69)] lie in the following arities:

\begin{tabular}{l|c|c|c}
entry & Weyl weight and degree of the pair & of the output & arity\\
\hline
$b\,c$ [BK, (3.64)] & $(2,1)$ & $(3,0)$ & $2$\\
$b\,b$ [BK, (3.65)] & $(4,0)$ & $(\tfrac92,0)$ & $1$\\
$\beta\,\gamma$ [BK, (3.66)] & $(\tfrac52,0)$ & $(3,0)$ & $1$\\
$\beta\,\beta$ [BK, (3.67)] & $(3,0)$ & $(\tfrac72,0)$ & $1$\\
$b\,\gamma$ [BK, (3.68)] & $(3,0)$ & $(\tfrac72,0)$ & $1$\\
$\beta\,b$ [BK, (3.69)] & $(\tfrac72,0)$ & $(4,0)$ & $1$\\
\end{tabular}
\tabcaption{The arities of the one-loop pair entries of the holomorphic twist, with $b$ represented by the self-dual field strength; with $b$ represented by a component of $\lambda^*$ the entry $b\,c$ has arity $1$.\label{tab:c11-twist}}

The entry $\beta\gamma$, the Konishi entry, is the two-body kernel of $Q_1$ along $\xi$ on the pair $(\lambda|_{S_+\otimes W},A|_{W_{\mathbb R}^\perp})$ with two outputs $\lambda|_{S_-\otimes W^*}$: the coefficient of $\int\omega\,(\bar\lambda_{\dot\alpha i}z^i,\bar\lambda^{\dot\alpha}{}_jz^j)$ in $Q_1$ of $\int\omega\,(u'^\alpha\lambda_\alpha{}^i,A_{ij}z^j)$. The classical component of this pair along $\xi$ is a number $\kappa$ times a cubic functional of the scalars, and a finite part of the operators of Weyl weight $3$ that adds $y\int\omega\,(\bar\lambda_{\dot\alpha i}z^i,\bar\lambda^{\dot\alpha}{}_jz^j)$ to that cubic functional changes the entry by $-\kappa y$ (Proposition \ref{prop:c11-freedom} (iv)), the two functionals having the same Weyl weight and quantum numbers. The pair entries are therefore compared through the class $[Q^\xi_1]$ on the cohomology of $Q^\xi_0$ (Theorem \ref{thm:c8-twist}), transported to the holomorphic letters through [ESW, Thm 10.9] (remark after Theorem \ref{thm:c8-twist}). Proposition \ref{prop:c11-plane} is the exact untwisted form, for two letters, of [BK, (3.77)], $Q_1(\gamma^I\gamma^J)=0$: for $\rho(w_a)\xi=0$, $\iota_{w_a}A$ is a combination of the components $A_{ij}z^j$, the letters $\gamma$. Theorem \ref{thm:c11-twistentries} gives the arity-one entries at $\psi=0$, and Theorem \ref{thm:c11-twistcheck} states what their comparison with [BK] tests.
\end{remark}

\begin{theorem}[The pair entries along a rank-one supercharge]
\label{thm:c11-twistentries}
Let $\xi=u\otimes z\in S_+\otimes W$ be of rank one, with $u'$, the twisted letters and their representatives of Remark \ref{rem:c11-twist}, and let $\psi=0$. Normalise the letters as the components along the weight vectors of a Cartan subalgebra of $\mathfrak{so}(V)\oplus\mathfrak{so}(W_{\mathbb R}^\perp)$ for which $u$, $u'$ and $z$ are weight vectors, read the outputs on the dual weight vectors, take $b$ as the component $F(d_1,d_2)$ of the field strength on the two null vectors $d_1$, $d_2$ of $V\otimes\mathbb C$ along which $\partial_{\dot\alpha}$ differentiates, and expand the output momenta in the null coordinates; this fixes each entry up to one nonzero factor for each letter and each output. Reduce modulo $\operatorname{Im}Q^\xi_0$, total derivatives kept: the derivatives $u^\alpha\sigma^m_{\alpha\dot\alpha}\partial_m$ are $Q^\xi_0$-exact on the letters, and so is, up to terms with antifields, the holomorphic Dirac component of a derivative on an output $\partial_{\dot\alpha}c$. Then the arity-one pair entries of $Q_1$ are those of Table \ref{tab:c11-twistentries}, and:
\begin{itemize}
\item (i) the entry $\beta\gamma$, the kernel of $(\lambda|_{S_+\otimes W},A|_{W_{\mathbb R}^\perp})$ on the leg set $(\lambda|_{S_-\otimes W^*})^2$, vanishes for every $\xi\in S_+\otimes W$, the placements on a line and at the special edge cancelling chain by chain; the one-body terms and $L_1^{(0)}$ have outputs of other Weyl weights or with a spectator letter, so the Konishi entry of Remark \ref{rem:c11-twist} is $0$ at $\psi=0$; the entry $bc$ with $b$ represented by $(\xi,\lambda^*)$ and the entry $\gamma\gamma$ vanish as well;
\item (ii) the entries $\beta\beta$ and $\beta b\to\gamma\gamma$ carry the colour and derivative structure of [BK, (3.67)] and of the last term of [BK, (3.69)]: the output $\gamma^K$, resp. the outputs $\gamma^J$, $\gamma^K$, with the indices complementary to those of the letters $\beta$, the same value in the two colour orders, and, for $\beta b\to\gamma\gamma$, the antisymmetric combination of the two holomorphic derivatives;
\item (iii) the entry $b\gamma$ has different values in the two colour orders, each with a term in $\log\tfrac32$, and the entry $\beta b\to\beta\,\partial c$ is nonzero in one colour order;
\item (iv) every derivative in the entries is holomorphic.
\end{itemize}
Statements (i)--(iv) are independent of the normalisation factors; the values of the individual entries depend on the scheme (Proposition \ref{prop:c11-freedom} (iv), Remark \ref{rem:c11-twist}).
\end{theorem}

\begin{tabular}{l|l|l|l|l}
entry & pair $\to$ outputs & order $(Y_1,Y_2)$ & order $(Y_2,Y_1)$ & Dirac components\\
\hline
$\beta\gamma$ [BK, (3.66)] & $(\beta_I,\gamma^J)\to\partial c\,\partial c$ & $0$ & $0$ & $0$, $0$\\
$bc$ [BK, (3.64)], $b\leftrightarrow(\xi,\lambda^*)$ & $(c,b)\to\partial c\,\partial c$ & $0$ & $0$ & $0$, $0$\\
$\beta\beta$ [BK, (3.67)] & $(\beta_I,\beta_J)\to\gamma^K\,\partial c$ & $\tfrac23$ & $\tfrac23$ & $\tfrac13-2L$, $\tfrac13-2L$\\
$b\gamma$ [BK, (3.68)] & $(b,\gamma^I)\to\gamma^I\,\partial c$ & $\tfrac{64}3-128L$ & $32+64L$ & $-\tfrac{64}3-128L$, $-\tfrac{32}3+64L$\\
$\beta b\to\beta\,\partial c$ [BK, (3.69)] & $(\beta_I,b)\to\beta_I\,\partial c$ & $-\tfrac{64}3$ & $0$ & $0$, $0$\\
$\beta b\to\gamma\gamma$ [BK, (3.69)] & $(\beta_I,b)\to\gamma^J\gamma^K$ & $\tfrac83\big(\partial_2\gamma^J\partial_1\gamma^K-\partial_1\gamma^J\partial_2\gamma^K\big)$ & the same & \\
$\gamma\gamma$ [BK, (3.77)] & $(\gamma^I,\gamma^J)\to$ twisted letters & $0$ & $0$ & \\
\end{tabular}
\tabcaption{The arity-one pair entries of $Q_1$ along $\xi=u\otimes z$ at $\psi=0$, reduced modulo $\operatorname{Im}Q^\xi_0$, in the normalisation of Theorem \ref{thm:c11-twistentries}: coefficients of $\frac{\hbar}{16\pi^2}$ times the colour factor $(1\otimes\mathrm{ad}_{Y_2}\mathrm{ad}_{Y_1})(-,-)^{-1}$ of the colour order $(Y_1,Y_2)$, $Y_1$ the first output listed, with $L=\log\tfrac32$. The third and fourth columns give the coefficient of the holomorphic derivative on the output other than $\partial c$, and for $\beta b\to\gamma\gamma$ the derivative structure, $\partial_1$ and $\partial_2$ being the derivatives along $d_1$ and $d_2$; the last column gives, in the two orders, the coefficient of the holomorphic Dirac component on the output $\partial c$, which is $Q^\xi_0$-exact up to terms with antifields; $K$ is complementary to $I,J$, and $J,K$ to $I$.\label{tab:c11-twistentries}}

\begin{proof}
The entries are the kernels of the pairs of Table \ref{tab:c11-defectpairs} on the leg sets $(\lambda|_{S_-\otimes W^*})^2$ of $(\lambda|_{S_+\otimes W},A|_{W_{\mathbb R}^\perp})$, $A|_{W_{\mathbb R}^\perp}\lambda|_{S_-\otimes W^*}$ of $(\lambda|_{S_+\otimes W},\lambda|_{S_+\otimes W})$ and, with one derivative, of $(A|_V,A|_{W_{\mathbb R}^\perp})$, $\lambda|_{S_+\otimes W}\lambda|_{S_-\otimes W^*}$ and $(A|_{W_{\mathbb R}^\perp})^2$ of $(\lambda|_{S_+\otimes W},A|_V)$ with one derivative, and $(\lambda|_{S_-\otimes W^*})^2$ of $(c,\lambda^*|_{S_+\otimes W^*})$; all have superficial degree $0$, so they carry no poles and no finite special vertex (Proposition \ref{prop:c11-oddbound}), and Theorem \ref{thm:c11-oddsimplex} evaluates them exactly, with the letter vectors contracted exactly in the weight basis (Remark \ref{rem:c11-oddmethod}). The colour factors are those of Lemma \ref{lem:c10-colour} with the output of a term $T_5$ read as an output of a cubic vertex. (i) The vanishing of $\beta\gamma$ for general $\xi$ and of $bc$ holds in both evaluations. The outputs $(\lambda|_{S_-\otimes W^*})^2$ have Weyl weight $3$, while $L_1^{(0)}$ preserves the Weyl weight $\tfrac52$ of the Konishi descendant, a one-body term keeps the other letter as a spectator, and the arity-two kernels have outputs of degree $-1$; hence the coefficient of $\int\omega\,(\bar\lambda_{\dot\alpha i}z^i,\bar\lambda^{\dot\alpha}{}_jz^j)$ in $\hbar Q^\xi_1$ of the Konishi descendant is the entry $\beta\gamma$. $\gamma\gamma$ is Proposition \ref{prop:c11-plane}. (ii)--(iv) are read off the table. Rescaling a weight vector rescales its letter by the inverse factor and its output by the factor, which multiplies an entry by one product of factors, the same for the two colour orders, for the reduced and the Dirac coefficients, and for the rational part and the part in $\log\tfrac32$; the vanishing, the equality of the two orders, the indices and the derivative structure are therefore independent of the normalisation.
\end{proof}

\begin{theorem}[What the comparison with the holomorphic twist tests]
\label{thm:c11-twistcheck}
Let $\xi=u\otimes z$ be of rank one, let $\Phi$ be the quasi-isomorphism from the cohomology of $Q^\xi_0$ on local operators to that of the free holomorphic differential on the letters $c$, $\gamma$, $\beta$, $b$, $\partial_{\dot\alpha}$ of [BK], transported through [ESW, Thm 10.9] (remark after Theorem \ref{thm:c8-twist}), and let $D$ be the difference, read through $\Phi$, of the one-loop quantisation $r_\xi(\hbar F_1)$ of the twisted theory restricted from the equivariant one and the holomorphic quantisation of [BK, (2.42), footnote 22] and [W18, Thm 3.4] (B. R. Williams, Renormalization for holomorphic field theories, arXiv:1809.02661). If $D\in\mathbb C\,S^\xi+\operatorname{Im}Q^\xi_0$, then $\Phi\circ[Q^\xi_1]\circ\Phi^{-1}$ is the class of the one-loop differential of [BK, (3.7), (3.64)--(3.69)], and:
\begin{itemize}
\item (a) the one-body terms of $Q^\xi_1$ on the twisted letters are $a\,Q^\xi_0+[Q^\xi_0,N]$ for a number $a$ and a derivation $N$ of the letters, zero on cohomology;
\item (b) on the pairs from $\{c,\gamma\}$ and on $(\beta,c)$ the two-body kernels lie, on the twisted letters, in $\operatorname{Im}Q^\xi_0$ plus terms cancelled on cohomology by the one-body terms, as in [BK, (3.11)]; on $(\gamma,\gamma)$ they vanish (Proposition \ref{prop:c11-plane});
\item (c) on the pairs $\beta\gamma$, $\beta\beta$, $b\gamma$, $\beta b$, $bb$ and $bc$ the kernels, reduced modulo $\operatorname{Im}Q^\xi_0$ as in Theorem \ref{thm:c11-twistentries}, have the colour factors $(1\otimes\mathrm{ad}_{Y_2}\mathrm{ad}_{Y_1})(-,-)^{-1}$, which are $\tfrac12f_{ACD}f_{BCE}$ of [BK] up to normalisation, and holomorphic derivatives $\partial_{\dot\alpha}=u'^\alpha\sigma^m_{\alpha\dot\alpha}\partial_m$ on the outputs;
\item (d) their coefficients satisfy, modulo the finite parts of Proposition \ref{prop:c11-freedom} (iv), the relations [BK, (3.59)--(3.63)], which are the order-$\hbar$ part of Corollary \ref{cor:c8-homotopy} (ii) at $(\xi,R)$ for the elements $R$ of the tree-level commutant $\mathfrak{psl}(3|3)$ of $\xi$, restricted by $r_\xi$;
\item (e) every statement of the holomorphic twist about $[Q_1]$ on classes of several letters is the image under $\Phi$ of $[Q^\xi_1]$, computed by the pair rule from the kernels of (c), with the normalisation of $\hbar$ relative to the unit fixed by $\Phi$.
\end{itemize}
For a general $D$ the two class maps differ by $[\{D,-\}]$, and (a)--(e) compute that class.
\end{theorem}

\begin{proof}
Two quantisations of the twisted theory differing by a $Q^\xi_0$-closed $D$ of degree $0$ have class maps differing by $[\{D,-\}]$, which vanishes for $D\in\mathbb C\,S^\xi+\operatorname{Im}Q^\xi_0$, since $\{S^\xi,-\}=Q^\xi_0$ and $\{Q^\xi_0\chi,-\}=\pm[Q^\xi_0,\{\chi,-\}]$ (Theorem \ref{thm:c8-twist} and the remark after it). (a) The one-body terms are brackets with local functionals (Theorem \ref{thm:c11-classes} (iii)), and the condition on $D$ makes them $aQ^\xi_0$ plus commutators with $Q^\xi_0$ on the twisted letters. (b), (c) The pair rule [BK, (3.7)--(3.11)], the colour factors of Lemma \ref{lem:c10-colour} with two outputs, and the exactness of the antiholomorphic derivatives, $\{Q,\bar Q^1_{\dot\alpha}\}=P_{-\dot\alpha}$ [BK, (2.5)]. (d) Corollary \ref{cor:c8-homotopy} (ii) at $(\xi,R)$ with $[\xi,R]=0$, at order $\hbar$, mapped by $r_\xi$ (Lemma \ref{lem:c8-ev}). (e) $[Q^\xi_1]$ is determined by its values on representatives, polynomials in the letters, through the pair rule.
\end{proof}

\begin{remark}[Method of computation]
\label{rem:c11-oddmethod}
Each defect cycle is evaluated exactly as in Remark \ref{rem:c11-method}: the Gaussian loop integral and the integrals of \eqref{eq:c11-oddsimplex} over the ordered subsimplices of the region where $\alpha_\sigma$ is the smallest coordinate, the letter derivatives through Proposition \ref{prop:c11-genfun} with $n+1$ edges, the position along $\eta$ and $K_b$ through the injected momentum $r$, the fibre algebra in an explicit realisation of $V_{10}$ and $\Sigma$ over $\mathbb Q(i)$ with the vertices $S^{(1)}(x)$ of Chapter 2 and the sign of Definition \ref{def:c11-defect} (ii), letter vectors with fractional components contracted exactly, and the colour words by Lemma \ref{lem:c10-colour}; the twisted letters are built from the weights of a Cartan subalgebra. A second evaluation works in momentum space with an independently written realisation of $S^{(1)}(x)$ in the components of Definition \ref{def:c8-twocomp}, the special edge as the factor $e^{-\varepsilon|p|^2}-1$, Feynman parameters for the other edges and numerical integration at several values of $\varepsilon$ with thirty digits, extrapolated to $\varepsilon^0$ with the poles fitted; along $K_b$ it evaluates the cycles at the endpoints $T=1,4,16$. The two evaluations give the same operators $C_\xi$, $Y_\eta$, $C_K$ and the same arrows of Lemma \ref{lem:c11-lines}, the same composite line along $\xi$, the same free directions of Proposition \ref{prop:c11-freexi}, the same kernels $\Upsilon_\xi$ on every leg set with at most two outputs of the pairs of field and ghost letters without derivatives and with one letter derivative, the same coefficients of Theorem \ref{thm:c11-bilinears} for both chiralities, the same entries $\beta\gamma$, $\beta\beta$ and $bc$, and the same translation-invariant kernels along $\eta$ and $K_b$ where both are evaluated; both reproduce Example \ref{ex:c11-cubic}, Corollary \ref{cor:c11-DPdefect} and Proposition \ref{prop:c11-plane}. The first evaluation extends $\Upsilon_\xi$ to three letter derivatives without the finite special vertices, and gives the entries of Table \ref{tab:c11-twistentries} and the identity (iv) of Theorem \ref{thm:c11-bilinears}; the second gives the closed forms of Theorem \ref{thm:c11-oddsimplex} (iv) for $n_t=2,\dots,6$, the finite special vertices on the pairs of gaugino letters, the pole identity \eqref{eq:c11-pole} on these pairs and on the pair of Example \ref{ex:c11-cubic} with the vacuum-pole terms, the position identity of Theorem \ref{thm:c11-oddsimplex} (ii), and, along $K_b$, $T\partial_T$ of the defect cycles of the intrinsic vertex $S^{(1)}(K_b)$ on $(A|_V,A|_V)$ with the output $A|_V$, with the vertex on a line and at the special edge, decreasing at least as $\varepsilon$ at the three endpoints, as Theorem \ref{thm:c11-oddsimplex} (iii) states. Every integral is a four-dimensional Schwinger integral with the cutoff $\varepsilon$ on every edge.
\end{remark}

\begin{remark}[The data that complete the one-loop components]
\label{rem:c11-next}
The one-body kernels of $Q_1$, $S_1$ and $K_1$ are given by the action side of the arity-one equations. Along $\xi$ it is \eqref{eq:c11-arity1xi}: the quadratic part of $\hbar\,\mathrm{Ob}^{(1)}(\xi)$ in the sectors $(\lambda,A|_V)$, $(\lambda,A|_{W_{\mathbb R}^\perp})$ and $(c,\lambda^*)$, the bracket $\{S^{(1)}(\xi),\hbar F_1^{(0)}\}$ with the finite counterterm of Theorem \ref{thm:c4-finite}, a particular solution, whose coefficient of $T_2$ gives $t_\perp$ and with it $c_Q$ of Theorem \ref{thm:c11-bilinears}, and the cubic part, which fixes $\theta^\xi_V-\theta^\xi_\perp$; the same computation along $\eta$ and $K_b$ gives the remaining one-body kernels. The next computations of the two-body kernels are the finite special vertices on the leg sets with $D\ge2$, with their vacuum-pole terms; the arity-two kernels along $(\xi,\xi)$, which give the entry $bc$ with $b$ represented by the self-dual field strength; and the entry $bb$ on the pair $(A|_V,A|_V)$ with two derivatives. With the one-body kernels, the relations (R2)--(R6) of Proposition \ref{prop:c11-relations} become finite computations on the scalar bilinears, with the cocycle finite part of Theorem \ref{thm:c11-bilinears} (iii) extended beyond the two-field sector; (R3) on the trace bilinear is the Konishi relation between $Q_1$, $S_1$ and $D_1=-48$ units. The class-level comparison of Theorem \ref{thm:c11-twistcheck} uses the transport map $\Phi$ built from [ESW, Thm 10.9], the condition on $D$ and the relations [BK, (3.59)--(3.63)] applied to the entries of Table \ref{tab:c11-twistentries}. The inner-product relations of Proposition \ref{prop:c11-S} (iv) and Corollary \ref{cor:c11-K} use $\mathrm{Gram}_{\rm 1}$ from the two-point graphs of Definition \ref{def:c8-2pt}, and $K_1\cong-(b,x)D_1$ on the families uses the rotation component $M_1$.
\end{remark}

\begin{remark}[Two-component forms]
\label{rem:c11-twocomp}
With the components of Definition \ref{def:c8-twocomp}:
\begin{itemize}
\item Letters. $A|_V$, $A|_{W_{\mathbb R}^\perp}$, $\lambda|_{S_+\otimes W}$, $\lambda|_{S_-\otimes W^*}$ are $A_m$, $A_{ij}$ (the $\mathbf 6$ of $SU(4)_R$), $\lambda_\alpha{}^i$ (undotted, $\mathbf 4$), $\bar\lambda_{\dot\alpha i}$ (dotted, $\bar{\mathbf 4}$), with the antifields $A^*_m$, $A^*_{ij}$, $\lambda^{*\alpha}{}_i$, $\bar\lambda^{*\dot\alpha i}$, $c^*$ and the degrees $1,0,0,0,0,-1,-1,-1,-1,-2$ for $c,A_m,A_{ij},\lambda_\alpha{}^i,\bar\lambda_{\dot\alpha i},A^*_m,A^*_{ij},\lambda^{*\alpha}{}_i,\bar\lambda^{*\dot\alpha i},c^*$; $(\lambda,\lambda^*)=(\lambda_\alpha{}^i,\lambda^{*\alpha}{}_i)-(\bar\lambda_{\dot\alpha i},\bar\lambda^{*\dot\alpha i})$ and $(u,w)=\tfrac14u^{ij}w_{ij}$ on $W_{\mathbb R}^\perp$.
\item Composite kernels. The internal vertex along $\Delta$ is $\int\big((2E+2)A_m,A^*_m\big)=2\int x_n(\partial_nA_m,A^*_m)+2\int(A_m,A^*_m)$ with $E=x_n\partial_n$, along $P_b$ it is $\int(b_n\partial_nA_m,A^*_m)$; it sits between $\langle A_mA_n\rangle=2(-,-)^{-1}\delta_{mn}$ and $\langle A^*_m\,c\rangle=-(-,-)^{-1}ip_m$, each times $e^{-\varepsilon|p|^2}|p|^{-2}$ at $T=\infty$, and the insertion factors are $2$, $2i(l+q_c)_n$ and $ib_n(l+q_c)_n$. The special edge carries $e^{-\varepsilon|p|^2}-1$ times $(-,-)^{-1}$ and the inverse fibre pairing of $(A_m,A^*_m)$, $\tfrac14(A^{ij},A^*_{ij})$, $(\lambda,\lambda^*)$ and $(c,c^*)$: on $(\lambda_\alpha{}^i,\lambda^{*\beta}{}_j)$ the factor $-|l+q_\sigma|^2\delta_\alpha^\beta\delta^i_j(-,-)^{-1}$, on $(\bar\lambda_{\dot\alpha i},\bar\lambda^{*\dot\beta j})$ the opposite sign.
\item Bilinears. For $O=\tfrac1{16}\int\omega\,(w_1^{ij}A_{ij},w_2^{kl}A_{kl})$ the composite pair kernel along $\Delta$ is
\[
\Big[-\tfrac13\big((2E+2)A_{ij}\big)A_{kl}-\tfrac13A_{ij}\big((2E+2)A_{kl}\big)-\tfrac43A_{ij}A_{kl}\Big]\ \text{units},
\]
and along $P_b$ it is $-\tfrac13$ units times $b_m\partial_m(A_{ij}A_{kl})$, with the same contraction.
\item Two-site density. On adjacent pairs $(A_{ij},A_{kl})$, $(Z,\psi)$, $(\psi,Z)$, $(\psi,\psi)$ of a single trace the composite pair kernels give $-\tfrac43$ units (translation-invariant part along $\Delta$) and $-\tfrac16$ units times $b_m\partial_m$ (along $P_b$) on the same word, and $0$ on the exchanged word.
\item Free directions. The sector parts of $S^{(1)}(\Delta)$ are $\tfrac14\int\big((2E+2)A^{ij},A^*_{ij}\big)$, $-\int\big[\big((2E+3)\lambda_\alpha{}^i,\lambda^{*\alpha}{}_i\big)-\big((2E+3)\bar\lambda_{\dot\alpha i},\bar\lambda^{*\dot\alpha i}\big)\big]$ and $\int\big((2E+2)A_m,A^*_m\big)-\int\big(2E(c),c^*\big)$; those of $S^{(1)}(P_b)$ are $\tfrac14\int(b_m\partial_mA^{ij},A^*_{ij})$, $-\int\big[(b_m\partial_m\lambda_\alpha{}^i,\lambda^{*\alpha}{}_i)-(b_m\partial_m\bar\lambda_{\dot\alpha i},\bar\lambda^{*\dot\alpha i})\big]$ and $\int(b_m\partial_mA_n,A^*_n)-\int(b_m\partial_mc,c^*)$. The exact directions are $Q\int x_m(A_m,c^*)$ and $Q\int(\rho(x)\lambda^*,\lambda^*)=-2i\,Q\int x_m\big(\lambda^{*\alpha}{}_i,\sigma^m_{\alpha\dot\alpha}\bar\lambda^{*\dot\alpha i}\big)$. The further class of Proposition \ref{prop:c11-tc0} without the reflection is a multiple of $\int\big[(\lambda_\alpha{}^i,\lambda^{*\alpha}{}_i)+(\bar\lambda_{\dot\alpha i},\bar\lambda^{*\dot\alpha i})\big]$, the rotation $\lambda_\alpha{}^i\mapsto e^{i\varphi}\lambda_\alpha{}^i$, $\bar\lambda_{\dot\alpha i}\mapsto e^{-i\varphi}\bar\lambda_{\dot\alpha i}$ of the free gaugino.
\item Letter count. It counts $2$ for each letter $A_m$, $A_{ij}$, $\lambda_\alpha{}^i$, $\bar\lambda_{\dot\alpha i}$ and $1$ for each derivative $\partial_m$ on it; (b) reads: the composite pair kernels along $P_b$ sum to $-\tfrac13$ units times $b_m\partial_m$ on cohomology.
\item Families and spectra. $n_m=\sigma^m_{1\dot1}$, $D_n=n_m(\partial_m+[A_m,-])$, $\psi=\lambda_1{}^i$, and $F_{11}$ the self-dual field strength with both undotted indices $1$. In Srednicki's normalisation, $A_m=-igA^a_mT^a$, $A_I=-igA^a_IT^a$, $\lambda=-\tfrac{ig}{\sqrt2}\lambda^aT^a$, $\operatorname{tr}(T^aT^b)=\tfrac12\delta^{ab}$, the unit is $-\frac{g^2N}{32\pi^2}$ and $-\tfrac13$ units is $\frac{g^2N}{96\pi^2}$. On the families $D_1=-8$ units times the eigenvalues $E$ of $\sum_lH_{l,l+1}$, which is $\frac{g^2N}{4\pi^2}E=2\gamma$ with the one-loop anomalous dimension $\gamma=\frac{g^2N}{8\pi^2}E$: $\gamma=\frac{g^2N}{2\pi^2}h(S)$ on twist two, $\gamma=\frac{3g^2N}{4\pi^2}$ for $\operatorname{tr}(F_{11}F_{11})$ and the trace bilinear, the Konishi multiplet.
\item Odd generators. $Q_1(O)=\xi_\alpha{}^i\,Q_1(O)^\alpha{}_i+\bar\xi_{\dot\alpha i}\,\bar Q_1(O)^{\dot\alpha i}$ and $S_1(O)=\eta^\alpha{}_i\,S_1(O)_\alpha{}^i+\bar\eta^{\dot\alpha i}\,\bar S_1(O)_{\dot\alpha i}$ (Definition \ref{def:c8-components}); the Killing spinor of $\eta$ has the components $(\rho(x)\eta)_\alpha{}^i=-i\,x_m\sigma^m_{\alpha\dot\alpha}\bar\eta^{\dot\alpha i}$ and $(\rho(x)\eta)_{\dot\alpha i}=i\,x_m\eta^\alpha{}_i\sigma^m_{\alpha\dot\alpha}$ \eqref{eq:c8-pairing}, so on a family with $P_1\cong0$ the position part along $b$ of $S_1$ along $\eta$ is, on cohomology, $Q_1$ along $\xi_\alpha{}^i=-i\,b_m\sigma^m_{\alpha\dot\alpha}\bar\eta^{\dot\alpha i}$, $\bar\xi_{\dot\alpha i}=i\,b_m\eta^\alpha{}_i\sigma^m_{\alpha\dot\alpha}$. For $\xi_\alpha{}^i=u_\alpha z^i$, $(\rho(w)\xi)^\alpha{}_j=\varepsilon^{\alpha\beta}w_{ji}u_\beta z^i$, so $\rho(w)\xi=0$ exactly when $w_{ij}z^j=0$, and then $\iota_wA=\tfrac14w^{ij}A_{ij}$ is a combination of the components $A_{ij}z^j$: Proposition \ref{prop:c11-plane} is the vanishing of $Q_1$ on $\int\omega\,(m_1^iA_{ij}z^j,m_2^kA_{kl}z^l)$, $m_1,m_2\in W$.
\item Terms of $S^{(1)}(\xi)$. For $\xi=(\xi_\alpha{}^i,0)$, with $(\xi^j\lambda^i):=\xi^{\alpha j}\lambda_\alpha{}^i$ and $F^0_{mn}=\partial_mA_n-\partial_nA_m$: $T_1=-i\big(\bar\lambda_{\dot\alpha i}\bar\sigma^{m\dot\alpha\alpha}\xi_\alpha{}^i,A^*_m\big)$, $T_2=\tfrac12\varepsilon_{jikl}\big((\xi^j\lambda^i),A^{*kl}\big)$, $T_3=-\tfrac14\big(F^0_{mn}(\sigma^m\bar\sigma^n)_\alpha{}^\beta\xi_\beta{}^i,\lambda^{*\alpha}{}_i\big)$, $T_4=-\tfrac i2\big(\xi^{\alpha j}\sigma^m_{\alpha\dot\alpha}\partial_mA_{ij},\bar\lambda^{*\dot\alpha i}\big)$, and $T_5$ the three terms $-\tfrac14\big([A_m,A_n](\sigma^m\bar\sigma^n)_\alpha{}^\beta\xi_\beta{}^i,\lambda^{*\alpha}{}_i\big)$, $-\tfrac i2\big(\xi^{\alpha j}\sigma^m_{\alpha\dot\alpha}[A_m,A_{ij}],\bar\lambda^{*\dot\alpha i}\big)$, $-\tfrac14\big([A^{ik},A_{kj}]\xi_\alpha{}^j,\lambda^{*\alpha}{}_i\big)$; for $\xi_\alpha{}^i=u_\alpha z^i$ these are the terms of the remark after Theorem \ref{thm:c8-twist}. The terms with $\bar\xi_{\dot\alpha i}$ follow by $\lambda_\alpha{}^i\leftrightarrow\bar\lambda_{\dot\alpha i}$, $\lambda^{*\alpha}{}_i\leftrightarrow\bar\lambda^{*\dot\alpha i}$, $\sigma\leftrightarrow\bar\sigma$. The intrinsic term along $\eta$ is $s_0=\big(A^{kj}\varepsilon_{\alpha\beta}\eta^\beta{}_j,\lambda^{*\alpha}{}_k\big)-\big(A_{kj}\varepsilon_{\dot\beta\dot\alpha}\bar\eta^{\dot\alpha j},\bar\lambda^{*\dot\beta k}\big)$, and $S^{(1)}(K_b)=\int\big\{(v.A_m,A^*_m)+\tfrac14(v.A^{ij},A^*_{ij})-(v.\lambda_\alpha{}^i,\lambda^{*\alpha}{}_i)+(v.\bar\lambda_{\dot\alpha i},\bar\lambda^{*\dot\alpha i})-(v.c,c^*)\big\}$ with $v^m=x^2b^m-2(b\cdot x)x^m$, $v.A_m=v^n\partial_nA_m+A_n\partial_mv^n$, $v.A_{ij}=v^n\partial_nA_{ij}-2(b\cdot x)A_{ij}$, $v.\lambda_\alpha{}^i=v^n\partial_n\lambda_\alpha{}^i-\tfrac12(x_mb_n-x_nb_m)(\sigma^m\bar\sigma^n)_\alpha{}^\beta\lambda_\beta{}^i-3(b\cdot x)\lambda_\alpha{}^i$, $v.c=v^n\partial_nc$.
\item Arrows and lines. $\xi.$ maps $A_{kl}\mapsto\Gamma_{kl}(\xi,\lambda)=-2\varepsilon_{jikl}(\xi^j\lambda^i)$ and $A_m\mapsto\Gamma^m(\xi,\lambda)=i\big(\bar\lambda_{\dot\alpha i}\bar\sigma^{m\dot\alpha\alpha}\xi_\alpha{}^i+\bar\xi_{\dot\alpha i}\bar\sigma^{m\dot\alpha\alpha}\lambda_\alpha{}^i\big)$. On a plane wave $C_\xi$ maps $\lambda$ to $c$ by $\lambda\mapsto-p_m\big(\bar\lambda_{\dot\alpha i}\bar\sigma^{m\dot\alpha\alpha}\xi_\alpha{}^i+\bar\xi_{\dot\alpha i}\bar\sigma^{m\dot\alpha\alpha}\lambda_\alpha{}^i\big)$. The composite line along $\xi_\alpha{}^i$ joins a leg $\lambda_\alpha{}^i$ to a leg $c$ through $\langle\lambda_\alpha{}^i\bar\lambda_{\dot\beta j}\rangle$, the term $T_1$ and $\langle A^*_m\,c\rangle$; with $\sigma^k_{\alpha\dot\beta}\bar\sigma^{m\dot\beta\gamma}p_kp_m=-|p|^2\delta_\alpha^\gamma$ it is $-2\,\xi_\alpha{}^i\,e^{-2\varepsilon|p|^2}|p|^{-2}(-,-)^{-1}$ on $c\otimes\lambda_\alpha{}^i$ and $+2\,\xi_\alpha{}^i\,e^{-2\varepsilon|p|^2}|p|^{-2}(-,-)^{-1}$ on $\lambda_\alpha{}^i\otimes c$. Along $\eta$, $(\rho(e_m)\eta)_\alpha{}^i=-i\sigma^m_{\alpha\dot\alpha}\bar\eta^{\dot\alpha i}$ and $(\rho(e_m)\eta)_{\dot\alpha i}=i\eta^\alpha{}_i\sigma^m_{\alpha\dot\alpha}$; along $K_b$, $\operatorname{div}v=-8\,b\cdot x$, $\Sigma^{\mu\nu}$ acts on $A_m$ as $(\Sigma^{\mu\nu}A)_m=\delta_{m\mu}A_\nu-\delta_{m\nu}A_\mu$ and on $\lambda_\alpha{}^i$ through $-\tfrac14(\sigma^\mu\bar\sigma^\nu-\sigma^\nu\bar\sigma^\mu)$, and $C_K$ maps $A_m$ to $-4b_mA_m$ in the slot of $c$.
\item Bilinears along $\xi$. For $O=\tfrac1{16}\int\omega\,(w_1^{ij}A_{ij},w_2^{kl}A_{kl})$,
\[
\mathcal B_1=\tfrac14\int\omega\,\big[(w_1^{kl}A_{kl},w_{2\,ji}\xi^{\alpha i}\lambda_\alpha{}^j)+(1\leftrightarrow2)\big],\qquad \mathcal B_2=\tfrac14w_1^{ij}w_{2\,ij}\int\omega\,(A_{ji},\xi^{\alpha i}\lambda_\alpha{}^j),
\]
$\mathcal B_3=\tfrac14w_1^{ij}w_{2\,ij}\int\omega\,i\big(\bar\lambda_{\dot\alpha i},A_m\bar\sigma^{m\dot\alpha\alpha}\xi_\alpha{}^i\big)$ and $\mathcal B_4=\tfrac14w_1^{ij}w_{2\,ij}\int\omega\,(c,\xi_\alpha{}^i\lambda^{*\alpha}{}_i)$; Theorem \ref{thm:c11-bilinears} reads, in units $-\frac{g^2N}{32\pi^2}$,
\[
Q_1(O)\big|_{\text{two-body}}=-2\log\tfrac32\,\mathcal B_1+4\log\tfrac32\,(\mathcal B_2+\mathcal B_3),\qquad L_1^{(0)}(O)=-\big(2+4\log\tfrac32\big)\tfrac14w_1^{ij}w_{2\,ij}\int\omega\,(\partial_mc,A_m),
\]
and the cocycle finite part is $\psi(O)=-\big(1+2\log\tfrac32\big)\tfrac14w_1^{ij}w_{2\,ij}\int\omega\,(A_m,A_m)$ units, $\psi(\mathcal B_1)=-2\,\mathcal B_3$ units. On the twenty traceless bilinears the transformation along $Q^i_\alpha$ is $1+c_Q$ times the classical one, with $c_Q$ of Theorem \ref{thm:c11-bilinears} (ii).
\item Special supercharges and special conformal fields on the bilinears. The two-body part of $S_1(O)$ along $\eta^\alpha{}_i$ is that of $\bar Q_1(O)$ along $\bar\xi_{\dot\alpha i}=i\,x_m\eta^\alpha{}_i\sigma^m_{\alpha\dot\alpha}$, and along $\bar\eta^{\dot\alpha i}$ that of $Q_1(O)$ along $\xi_\alpha{}^i=-i\,x_m\sigma^m_{\alpha\dot\alpha}\bar\eta^{\dot\alpha i}$. On the part with two scalar legs, $K_1(O)=-\big(\tfrac13+\theta\big)\big(v^m\partial_mO-4(b\cdot x)O\big)$ units on the traceless bilinears and the same plus $48\,(b\cdot x)\,O$ units on the trace $\tfrac14\int\omega\,(A^{ij},A_{ij})$.
\item Twist. With $z=e_4$: $\gamma^I=A_{I4}$, $\beta_I=u'^\alpha\lambda_\alpha{}^I$ ($I=1,2,3$), $\partial_{\dot\alpha}c=\bar\lambda_{\dot\alpha4}$, $b=u'^\alpha F_{mn}(\sigma^m\bar\sigma^n)_\alpha{}^\beta u'_\beta$ or $(\xi,\lambda^*)=u_\alpha\lambda^{*\alpha}{}_4$, and $\partial_{\dot\alpha}=u'^\alpha\sigma^m_{\alpha\dot\alpha}D_m$, so that $\partial_2\gamma^J\partial_1\gamma^K-\partial_1\gamma^J\partial_2\gamma^K$ of Table \ref{tab:c11-twistentries} is a multiple of $\partial_{\dot\alpha}\gamma^J\partial^{\dot\alpha}\gamma^K$. The Konishi entry is the coefficient of $\int\omega\,(\bar\lambda_{\dot\alpha4},\bar\lambda^{\dot\alpha}{}_4)$ in $Q_1\big(\int\omega\,(u'^\alpha\lambda_\alpha{}^i,A_{i4})\big)$, $0$ at $\psi=0$; [BK, (3.66)] is $Q_1\big((\beta_I)^A(\gamma^J)^B\big)=\tfrac12\delta^J_If_{ACD}f_{BCE}\,\partial_{\dot\alpha}c^D\partial^{\dot\alpha}c^E$, and the entries $\beta\beta$, $b\gamma$, $\beta b$ of [BK, (3.67)--(3.69)] have the same factor $\tfrac12f_{ACD}f_{BCE}$ in both terms.
\item Weinberg. On the Dirac column $(\lambda_\alpha{}^i,\bar\lambda^{\dot\alpha}{}_i)$, $\rho(v)=\gamma_W^\mu v_\mu$ gives Weinberg's form of each line, as in Remark \ref{rem:c10-twocomp}.
\end{itemize}
\end{remark}
